% generated by GAPDoc2LaTeX from XML source (Frank Luebeck)
\documentclass[a4paper,11pt]{report}

\usepackage[top=37mm,bottom=37mm,left=27mm,right=27mm]{geometry}
\sloppy
\pagestyle{myheadings}
\usepackage{amssymb}
\usepackage[utf8]{inputenc}
\usepackage{makeidx}
\makeindex
\usepackage{color}
\definecolor{FireBrick}{rgb}{0.5812,0.0074,0.0083}
\definecolor{RoyalBlue}{rgb}{0.0236,0.0894,0.6179}
\definecolor{RoyalGreen}{rgb}{0.0236,0.6179,0.0894}
\definecolor{RoyalRed}{rgb}{0.6179,0.0236,0.0894}
\definecolor{LightBlue}{rgb}{0.8544,0.9511,1.0000}
\definecolor{Black}{rgb}{0.0,0.0,0.0}

\definecolor{linkColor}{rgb}{0.0,0.0,0.554}
\definecolor{citeColor}{rgb}{0.0,0.0,0.554}
\definecolor{fileColor}{rgb}{0.0,0.0,0.554}
\definecolor{urlColor}{rgb}{0.0,0.0,0.554}
\definecolor{promptColor}{rgb}{0.0,0.0,0.589}
\definecolor{brkpromptColor}{rgb}{0.589,0.0,0.0}
\definecolor{gapinputColor}{rgb}{0.589,0.0,0.0}
\definecolor{gapoutputColor}{rgb}{0.0,0.0,0.0}

%%  for a long time these were red and blue by default,
%%  now black, but keep variables to overwrite
\definecolor{FuncColor}{rgb}{0.0,0.0,0.0}
%% strange name because of pdflatex bug:
\definecolor{Chapter }{rgb}{0.0,0.0,0.0}
\definecolor{DarkOlive}{rgb}{0.1047,0.2412,0.0064}


\usepackage{fancyvrb}

\usepackage{mathptmx,helvet}
\usepackage[T1]{fontenc}
\usepackage{textcomp}


\usepackage[
            pdftex=true,
            bookmarks=true,        
            a4paper=true,
            pdftitle={Written with GAPDoc},
            pdfcreator={LaTeX with hyperref package / GAPDoc},
            colorlinks=true,
            backref=page,
            breaklinks=true,
            linkcolor=linkColor,
            citecolor=citeColor,
            filecolor=fileColor,
            urlcolor=urlColor,
            pdfpagemode={UseNone}, 
           ]{hyperref}

\newcommand{\maintitlesize}{\fontsize{50}{55}\selectfont}

% write page numbers to a .pnr log file for online help
\newwrite\pagenrlog
\immediate\openout\pagenrlog =\jobname.pnr
\immediate\write\pagenrlog{PAGENRS := [}
\newcommand{\logpage}[1]{\protect\write\pagenrlog{#1, \thepage,}}
%% were never documented, give conflicts with some additional packages

\newcommand{\GAP}{\textsf{GAP}}

%% nicer description environments, allows long labels
\usepackage{enumitem}
\setdescription{style=nextline}

%% depth of toc
\setcounter{tocdepth}{1}



\usepackage{graphicx,amsfonts,amsmath, amssymb}

%% command for ColorPrompt style examples
\newcommand{\gapprompt}[1]{\color{promptColor}{\bfseries #1}}
\newcommand{\gapbrkprompt}[1]{\color{brkpromptColor}{\bfseries #1}}
\newcommand{\gapinput}[1]{\color{gapinputColor}{#1}}


\begin{document}

\logpage{[ 0, 0, 0 ]}
\begin{titlepage}
\mbox{}\vfill

\begin{center}{\maintitlesize \textbf{ HypArr \mbox{}}}\\
\vfill

\hypersetup{pdftitle= HypArr }
\markright{\scriptsize \mbox{}\hfill  HypArr  \hfill\mbox{}}
{\Huge \textbf{ Computations with hyperplane arrangements \mbox{}}}\\
\vfill

{\Huge  0.26.04 \mbox{}}\\[1cm]
{ April 2, 2026 \mbox{}}\\[1cm]
\mbox{}\\[2cm]
{\Large \textbf{ Paul M{\"u}cksch\\
    \mbox{}}}\\
\hypersetup{pdfauthor= Paul M{\"u}cksch\\
    }
\end{center}\vfill

\mbox{}\\
{\mbox{}\\
\small \noindent \textbf{ Paul M{\"u}cksch\\
    }  Email: \href{mailto://paul.muecksch+uni@gmail.com} {\texttt{paul.muecksch+uni@gmail.com}}\\
  Homepage: \href{https://paulmuecksch.github.io} {\texttt{https://paulmuecksch.github.io}}\\
  Address: \begin{minipage}[t]{8cm}\noindent
 Leibniz Universit{\"a}t Hannover\\
 Institut f{\"u}r Algebra, Zahlentheorie\\
 und Diskrete Mathematik\\
 Welfengarten 1\\
 30167 Hannover, Germany\\
 \end{minipage}
}\\
\end{titlepage}

\newpage\setcounter{page}{2}
\newpage

\def\contentsname{Contents\logpage{[ 0, 0, 1 ]}}

\tableofcontents
\newpage

     
\chapter{\textcolor{Chapter }{Introduction}}\label{Chapter_Introduction}
\logpage{[ 1, 0, 0 ]}
\hyperdef{L}{X7DFB63A97E67C0A1}{}
{
  

 \textsc{HypArr} is a \textsc{GAP 4} package for computations with hyperplane arrangements (see \cite{OrTer92_Arr}), oriented matroids (see \cite{BLSWZ1999_OrientedMatroids}), their combinatorial and topological invariants, in particular geometric
lattices, face posets, Milnor fibers and complements. 

 }

   
\chapter{\textcolor{Chapter }{Hyperplane arrangements}}\label{Chapter_Hyperplane_arrangements}
\logpage{[ 2, 0, 0 ]}
\hyperdef{L}{X7EF7AC44814961EC}{}
{
  

 Let $\mathbb{K}$ be a field. A hyperplane arrangment is a finite set $\mathcal{A} = \{H_1,\ldots,H_n\}$ of hyperplanes in a finite dimensional $\mathbb{K}$\texttt{\symbol{45}}vector space $V$. 
\section{\textcolor{Chapter }{Construction of arrangements of hyperplanes}}\label{Chapter_Hyperplane_arrangements_Section_Construction_of_arrangements_of_hyperplanes}
\logpage{[ 2, 1, 0 ]}
\hyperdef{L}{X80F9235385C15A41}{}
{
  

\subsection{\textcolor{Chapter }{HyperplaneArrangement (for IsList)}}
\logpage{[ 2, 1, 1 ]}\nobreak
\hyperdef{L}{X860548D886CE5BA2}{}
{\noindent\textcolor{FuncColor}{$\triangleright$\enspace\texttt{HyperplaneArrangement({\mdseries\slshape r})\index{HyperplaneArrangement@\texttt{HyperplaneArrangement}!for IsList}
\label{HyperplaneArrangement:for IsList}
}\hfill{\scriptsize (operation)}}\\
\textbf{\indent Returns:\ }
A hyperplane arrangement 



 Constructs a hyperplane arrangement from a list \mbox{\texttt{\mdseries\slshape R}} of vectors representing defining linear forms of hyperplanes. 

 Each vector $r = [r_1,\dots,r_n]$ corresponds to the hyperplane 
\[ r_1 x_1 + \cdots + r_n x_n = 0. \]
 The vectors must lie in the same vector space. 

 Linearly dependent defining forms are removed automatically, so the resulting
arrangement only stores pairwise linearly independent hyperplanes. 

 
\begin{Verbatim}[commandchars=!@|,fontsize=\small,frame=single,label=Example]
  !gapprompt@gap>| !gapinput@A := HyperplaneArrangement([[1,0,0],[0,1,0],[0,0,1]]);|
  <HyperplaneArrangement: 3 hyperplanes in 3-space>
  !gapprompt@gap>| !gapinput@Roots(A);|
  [ [ 1, 0, 0 ], [ 0, 1, 0 ], [ 0, 0, 1 ] ]
\end{Verbatim}
 }

 }

 
\section{\textcolor{Chapter }{Attributes of arrangements}}\label{Chapter_Hyperplane_arrangements_Section_Attributes_of_arrangements}
\logpage{[ 2, 2, 0 ]}
\hyperdef{L}{X8278B3F07C13FCE8}{}
{
  

\subsection{\textcolor{Chapter }{Dimension (for IsHyperplaneArrangement)}}
\logpage{[ 2, 2, 1 ]}\nobreak
\hyperdef{L}{X7CCC1A2A7C3F2CC5}{}
{\noindent\textcolor{FuncColor}{$\triangleright$\enspace\texttt{Dimension({\mdseries\slshape A})\index{Dimension@\texttt{Dimension}!for IsHyperplaneArrangement}
\label{Dimension:for IsHyperplaneArrangement}
}\hfill{\scriptsize (attribute)}}\\
\textbf{\indent Returns:\ }
A non\texttt{\symbol{45}}negative integer. 



 Returns the dimension of the ambient vector space in which the hyperplane
arrangement \mbox{\texttt{\mdseries\slshape A}} is defined. 
\begin{Verbatim}[commandchars=!@|,fontsize=\small,frame=single,label=Example]
  !gapprompt@gap>| !gapinput@A:=HyperplaneArrangement([[1,0,0],[0,1,0]]); Dimension(A);|
  <HyperplaneArrangement: 2 hyperplanes in 3-space>
  3
\end{Verbatim}
 }

 

\subsection{\textcolor{Chapter }{Roots (for IsHyperplaneArrangement)}}
\logpage{[ 2, 2, 2 ]}\nobreak
\hyperdef{L}{X8419FA5F7A08DF56}{}
{\noindent\textcolor{FuncColor}{$\triangleright$\enspace\texttt{Roots({\mdseries\slshape A})\index{Roots@\texttt{Roots}!for IsHyperplaneArrangement}
\label{Roots:for IsHyperplaneArrangement}
}\hfill{\scriptsize (attribute)}}\\
\textbf{\indent Returns:\ }
A list of vectors. 



 Returns the list of vectors giving defining linear forms (also called \emph{roots}) of the hyperplanes in the arrangement \mbox{\texttt{\mdseries\slshape A}}. 
\begin{Verbatim}[commandchars=!@|,fontsize=\small,frame=single,label=Example]
  !gapprompt@gap>| !gapinput@A:=AGpql(3,3,2); Roots(A);|
  <HyperplaneArrangement: 3 hyperplanes in 2-space>
  [ [ 1, -E(3) ], [ 1, -E(3)^2 ], [ 1, -1 ] ]
\end{Verbatim}
 }

 

\subsection{\textcolor{Chapter }{IntersectionLattice (for IsHyperplaneArrangement)}}
\logpage{[ 2, 2, 3 ]}\nobreak
\hyperdef{L}{X8555074B83058460}{}
{\noindent\textcolor{FuncColor}{$\triangleright$\enspace\texttt{IntersectionLattice({\mdseries\slshape A})\index{IntersectionLattice@\texttt{IntersectionLattice}!for IsHyperplaneArrangement}
\label{IntersectionLattice:for IsHyperplaneArrangement}
}\hfill{\scriptsize (attribute)}}\\
\textbf{\indent Returns:\ }
An object of type \texttt{IsGeomLattice}. 



 Computes the intersection lattice of the hyperplane arrangement \mbox{\texttt{\mdseries\slshape A}}. 

 The ground set \texttt{L} (see \texttt{LGroundSet} (\ref{LGroundSet:for IsGeomLattice})) of the lattice is represented
level\texttt{\symbol{45}}by\texttt{\symbol{45}}level. The entry \texttt{L[k]} contains all intersections of rank $k$. 

 Each lattice element is stored as a list of indices referring to hyperplanes
in \texttt{Roots(A)}. 

 For example, the list \texttt{[1,3]} represents the intersection of the first and third hyperplane of the
arrangement. 

 
\begin{Verbatim}[commandchars=!@|,fontsize=\small,frame=single,label=Example]
  !gapprompt@gap>| !gapinput@A := HyperplaneArrangement([[1,0,0],[0,1,0],[0,0,0],[1,1,1]]);|
  <HyperplaneArrangement: 3 hyperplanes in 3-space>
  !gapprompt@gap>| !gapinput@L:=IntersectionLattice(A); LGroundSet(L);|
  <Geometric lattice: 3 atoms, rank 3>
  [ [ [ 1 ], [ 2 ], [ 3 ] ], 
    [ [ 1, 2 ], [ 1, 3 ], [ 2, 3 ] ], 
    [ [ 1, 2, 3 ] ] ]
\end{Verbatim}
 }

 

\subsection{\textcolor{Chapter }{CharPoly (for IsHyperplaneArrangement)}}
\logpage{[ 2, 2, 4 ]}\nobreak
\hyperdef{L}{X8426CC4D7A8491A5}{}
{\noindent\textcolor{FuncColor}{$\triangleright$\enspace\texttt{CharPoly({\mdseries\slshape A})\index{CharPoly@\texttt{CharPoly}!for IsHyperplaneArrangement}
\label{CharPoly:for IsHyperplaneArrangement}
}\hfill{\scriptsize (attribute)}}\\
\textbf{\indent Returns:\ }
a univariate polynomial with integral coefficients. 



 Returns $\chi_{\mathcal{A}}$ the \emph{characterisitc polynomial} of the arrangement \mbox{\texttt{\mdseries\slshape A}}. 
\begin{Verbatim}[commandchars=!@|,fontsize=\small,frame=single,label=Example]
  !gapprompt@gap>| !gapinput@A:=AGpql(2,2,3);|
  <HyperplaneArrangement: 6 hyperplanes in 3-space>
  !gapprompt@gap>| !gapinput@CharPoly(A);|
  t^3-6*t^2+11*t-6
\end{Verbatim}
 }

 

\subsection{\textcolor{Chapter }{ExpArr (for IsHyperplaneArrangement)}}
\logpage{[ 2, 2, 5 ]}\nobreak
\hyperdef{L}{X83DFD7407D949CD9}{}
{\noindent\textcolor{FuncColor}{$\triangleright$\enspace\texttt{ExpArr({\mdseries\slshape A})\index{ExpArr@\texttt{ExpArr}!for IsHyperplaneArrangement}
\label{ExpArr:for IsHyperplaneArrangement}
}\hfill{\scriptsize (attribute)}}\\
\textbf{\indent Returns:\ }
a list of integers, or fail 



 Returns \emph{exponents} of the arrangement $\mathcal{A}$ if the characteristic polynomial factors over the integers. 
\begin{Verbatim}[commandchars=!@|,fontsize=\small,frame=single,label=Example]
  !gapprompt@gap>| !gapinput@A:=HyperplaneArrangement([[1,0],[0,1],[1,1]]);|
  <HyperplaneArrangement: 3 hyperplanes in 2-space>
  !gapprompt@gap>| !gapinput@ExpArr(A);|
  [ 1, 2 ]
\end{Verbatim}
 }

 

\subsection{\textcolor{Chapter }{MSetInvL (for IsHyperplaneArrangement)}}
\logpage{[ 2, 2, 6 ]}\nobreak
\hyperdef{L}{X852CC1A8795FCD15}{}
{\noindent\textcolor{FuncColor}{$\triangleright$\enspace\texttt{MSetInvL({\mdseries\slshape A})\index{MSetInvL@\texttt{MSetInvL}!for IsHyperplaneArrangement}
\label{MSetInvL:for IsHyperplaneArrangement}
}\hfill{\scriptsize (attribute)}}\\
\textbf{\indent Returns:\ }
An object of type IsGeomLattice. 



 Computes multiset invariants of the intersection lattice (see \texttt{IntersectionLattice} (\ref{IntersectionLattice:for IsHyperplaneArrangement})) of the hyperplane arrangement \mbox{\texttt{\mdseries\slshape A}}. 

 For each level of the lattice, the function records how many intersections
occur with a given number of defining hyperplanes. 

 
\begin{Verbatim}[commandchars=!@|,fontsize=\small,frame=single,label=Example]
  !gapprompt@gap>| !gapinput@A:=HyperplaneArrangement([[1,0,0,],[0,1,0],[0,0,1],[1,1,1]]);|
  <HyperplaneArrangement: 4 hyperplanes in 3-space>
  !gapprompt@gap>| !gapinput@MSetInvL(A);|
  [ [ [ 1, 4 ] ], [ [ 2, 6 ] ], [ [ 4, 1 ] ] ]
\end{Verbatim}
 }

 }

 
\section{\textcolor{Chapter }{Properties of arrangements}}\label{Chapter_Hyperplane_arrangements_Section_Properties_of_arrangements}
\logpage{[ 2, 3, 0 ]}
\hyperdef{L}{X7F6F9A097C0400D3}{}
{
  

\subsection{\textcolor{Chapter }{IsReal (for IsHyperplaneArrangement)}}
\logpage{[ 2, 3, 1 ]}\nobreak
\hyperdef{L}{X7E907A097A17195F}{}
{\noindent\textcolor{FuncColor}{$\triangleright$\enspace\texttt{IsReal({\mdseries\slshape A})\index{IsReal@\texttt{IsReal}!for IsHyperplaneArrangement}
\label{IsReal:for IsHyperplaneArrangement}
}\hfill{\scriptsize (property)}}\\
\textbf{\indent Returns:\ }
true or false. 



 Determines, if the hyperplane arrangement is defined over the reals, i.e., if
the entries of the defining linear forms are real. }

 

\subsection{\textcolor{Chapter }{HArrIsIrreducible (for IsHyperplaneArrangement)}}
\logpage{[ 2, 3, 2 ]}\nobreak
\hyperdef{L}{X85260BFF870DE234}{}
{\noindent\textcolor{FuncColor}{$\triangleright$\enspace\texttt{HArrIsIrreducible({\mdseries\slshape A})\index{HArrIsIrreducible@\texttt{HArrIsIrreducible}!for IsHyperplaneArrangement}
\label{HArrIsIrreducible:for IsHyperplaneArrangement}
}\hfill{\scriptsize (property)}}\\
\textbf{\indent Returns:\ }
true or false. 



 Determines, if the hyperplane arrangement is irreducible. }

 

\subsection{\textcolor{Chapter }{HArrIsGeneric (for IsHyperplaneArrangement)}}
\logpage{[ 2, 3, 3 ]}\nobreak
\hyperdef{L}{X804526007DD22027}{}
{\noindent\textcolor{FuncColor}{$\triangleright$\enspace\texttt{HArrIsGeneric({\mdseries\slshape A})\index{HArrIsGeneric@\texttt{HArrIsGeneric}!for IsHyperplaneArrangement}
\label{HArrIsGeneric:for IsHyperplaneArrangement}
}\hfill{\scriptsize (property)}}\\
\textbf{\indent Returns:\ }
true or false. 



 Determines, if the hyperplane arrangement is generic. See also \texttt{LIsGeneric} (\ref{LIsGeneric:for IsGeomLattice}). 
\begin{Verbatim}[commandchars=!@|,fontsize=\small,frame=single,label=Example]
  !gapprompt@gap>| !gapinput@A := HyperplaneArrangement([[1,0,0],[0,1,0],[0,0,1],[1,1,0],[1,1,1]]);|
  <HyperplaneArrangement: 5 hyperplanes in 3-space>
  !gapprompt@gap>| !gapinput@HArrIsGeneric(A);|
  false
  !gapprompt@gap>| !gapinput@A := HyperplaneArrangement([[1,0,0],[0,1,0],[0,0,1],[1,1,1]]);|
  <HyperplaneArrangement: 4 hyperplanes in 3-space>
  !gapprompt@gap>| !gapinput@HArrIsGeneric(A);|
  true
\end{Verbatim}
 }

 }

 
\section{\textcolor{Chapter }{Attributes of geometric lattices}}\label{Chapter_Hyperplane_arrangements_Section_Attributes_of_geometric_lattices}
\logpage{[ 2, 4, 0 ]}
\hyperdef{L}{X79DB21318677B969}{}
{
  

\subsection{\textcolor{Chapter }{LGroundSet (for IsGeomLattice)}}
\logpage{[ 2, 4, 1 ]}\nobreak
\hyperdef{L}{X864BB73F7812E0B4}{}
{\noindent\textcolor{FuncColor}{$\triangleright$\enspace\texttt{LGroundSet({\mdseries\slshape L})\index{LGroundSet@\texttt{LGroundSet}!for IsGeomLattice}
\label{LGroundSet:for IsGeomLattice}
}\hfill{\scriptsize (attribute)}}\\
\textbf{\indent Returns:\ }
list of lists 



 Returns the ground set of the geometric lattice \mbox{\texttt{\mdseries\slshape L}}. }

 

\subsection{\textcolor{Chapter }{LAtoms (for IsGeomLattice)}}
\logpage{[ 2, 4, 2 ]}\nobreak
\hyperdef{L}{X83946B69829FDA91}{}
{\noindent\textcolor{FuncColor}{$\triangleright$\enspace\texttt{LAtoms({\mdseries\slshape L})\index{LAtoms@\texttt{LAtoms}!for IsGeomLattice}
\label{LAtoms:for IsGeomLattice}
}\hfill{\scriptsize (attribute)}}\\
\textbf{\indent Returns:\ }
list 



 Returns the set of atoms of \mbox{\texttt{\mdseries\slshape L}}. }

 

\subsection{\textcolor{Chapter }{LRank (for IsGeomLattice)}}
\logpage{[ 2, 4, 3 ]}\nobreak
\hyperdef{L}{X86823DA184430088}{}
{\noindent\textcolor{FuncColor}{$\triangleright$\enspace\texttt{LRank({\mdseries\slshape L})\index{LRank@\texttt{LRank}!for IsGeomLattice}
\label{LRank:for IsGeomLattice}
}\hfill{\scriptsize (attribute)}}\\
\textbf{\indent Returns:\ }
a non\texttt{\symbol{45}}negative integer 



 Returns the rank of \mbox{\texttt{\mdseries\slshape L}}. }

 

\subsection{\textcolor{Chapter }{LkFlats (for IsGeomLattice)}}
\logpage{[ 2, 4, 4 ]}\nobreak
\hyperdef{L}{X8255DDFF86DBA263}{}
{\noindent\textcolor{FuncColor}{$\triangleright$\enspace\texttt{LkFlats({\mdseries\slshape L})\index{LkFlats@\texttt{LkFlats}!for IsGeomLattice}
\label{LkFlats:for IsGeomLattice}
}\hfill{\scriptsize (attribute)}}\\
\textbf{\indent Returns:\ }
a function 



 Returns a function extracting the flats of rank $k$ in \mbox{\texttt{\mdseries\slshape L}}. }

 

\subsection{\textcolor{Chapter }{LRankFunction (for IsGeomLattice)}}
\logpage{[ 2, 4, 5 ]}\nobreak
\hyperdef{L}{X807645CC7B8E3782}{}
{\noindent\textcolor{FuncColor}{$\triangleright$\enspace\texttt{LRankFunction({\mdseries\slshape L})\index{LRankFunction@\texttt{LRankFunction}!for IsGeomLattice}
\label{LRankFunction:for IsGeomLattice}
}\hfill{\scriptsize (attribute)}}\\
\textbf{\indent Returns:\ }
A function 



 Returns the rank function of the geometric lattice \mbox{\texttt{\mdseries\slshape L}}. }

 

\subsection{\textcolor{Chapter }{LGraph (for IsGeomLattice)}}
\logpage{[ 2, 4, 6 ]}\nobreak
\hyperdef{L}{X79A909F786A94A47}{}
{\noindent\textcolor{FuncColor}{$\triangleright$\enspace\texttt{LGraph({\mdseries\slshape L})\index{LGraph@\texttt{LGraph}!for IsGeomLattice}
\label{LGraph:for IsGeomLattice}
}\hfill{\scriptsize (attribute)}}\\
\textbf{\indent Returns:\ }
a graph 



 Constructs the directed graph of the Hasse diagram of L. }

 

\subsection{\textcolor{Chapter }{LAutGroup (for IsGeomLattice)}}
\logpage{[ 2, 4, 7 ]}\nobreak
\hyperdef{L}{X7E1203CB7969B4A0}{}
{\noindent\textcolor{FuncColor}{$\triangleright$\enspace\texttt{LAutGroup({\mdseries\slshape L})\index{LAutGroup@\texttt{LAutGroup}!for IsGeomLattice}
\label{LAutGroup:for IsGeomLattice}
}\hfill{\scriptsize (attribute)}}\\
\textbf{\indent Returns:\ }
a group 



 Computes the autmorphism group of L as a subgroup of \texttt{Sym(LAtoms(L))}. }

 

\subsection{\textcolor{Chapter }{LMoebius (for IsGeomLattice)}}
\logpage{[ 2, 4, 8 ]}\nobreak
\hyperdef{L}{X80A93EFA7EC0496A}{}
{\noindent\textcolor{FuncColor}{$\triangleright$\enspace\texttt{LMoebius({\mdseries\slshape L})\index{LMoebius@\texttt{LMoebius}!for IsGeomLattice}
\label{LMoebius:for IsGeomLattice}
}\hfill{\scriptsize (attribute)}}\\
\textbf{\indent Returns:\ }
function 



 Returns the Moebius function of \mbox{\texttt{\mdseries\slshape L}}. }

 

\subsection{\textcolor{Chapter }{LCharPoly (for IsGeomLattice)}}
\logpage{[ 2, 4, 9 ]}\nobreak
\hyperdef{L}{X83B3C4237FD0B9AA}{}
{\noindent\textcolor{FuncColor}{$\triangleright$\enspace\texttt{LCharPoly({\mdseries\slshape L})\index{LCharPoly@\texttt{LCharPoly}!for IsGeomLattice}
\label{LCharPoly:for IsGeomLattice}
}\hfill{\scriptsize (attribute)}}\\
\textbf{\indent Returns:\ }
function 



 Computes the characteristic polynomial of \mbox{\texttt{\mdseries\slshape L}}. }

 }

 
\section{\textcolor{Chapter }{Properties of geometric lattices}}\label{Chapter_Hyperplane_arrangements_Section_Properties_of_geometric_lattices}
\logpage{[ 2, 5, 0 ]}
\hyperdef{L}{X8056EC0986604552}{}
{
  

\subsection{\textcolor{Chapter }{LIsIrreducible (for IsGeomLattice)}}
\logpage{[ 2, 5, 1 ]}\nobreak
\hyperdef{L}{X85C9C1F47977DAD5}{}
{\noindent\textcolor{FuncColor}{$\triangleright$\enspace\texttt{LIsIrreducible({\mdseries\slshape L})\index{LIsIrreducible@\texttt{LIsIrreducible}!for IsGeomLattice}
\label{LIsIrreducible:for IsGeomLattice}
}\hfill{\scriptsize (property)}}\\
\textbf{\indent Returns:\ }
true or false. 



 Determines, if L is irreducible. }

 

\subsection{\textcolor{Chapter }{LIsBoolean (for IsGeomLattice)}}
\logpage{[ 2, 5, 2 ]}\nobreak
\hyperdef{L}{X8444E8837E952FE4}{}
{\noindent\textcolor{FuncColor}{$\triangleright$\enspace\texttt{LIsBoolean({\mdseries\slshape L})\index{LIsBoolean@\texttt{LIsBoolean}!for IsGeomLattice}
\label{LIsBoolean:for IsGeomLattice}
}\hfill{\scriptsize (property)}}\\
\textbf{\indent Returns:\ }
true or false. 



 Determines, if L is a boolean lattice. 
\begin{Verbatim}[commandchars=!@|,fontsize=\small,frame=single,label=Example]
  !gapprompt@gap>| !gapinput@L:=IntersectionLattice(HyperplaneArrangement([[1,0],[1,1]]));|
  <Geometric lattice: 2 atoms, rank 2>
  !gapprompt@gap>| !gapinput@LIsBoolean(L);|
  true
  !gapprompt@gap>| !gapinput@L:=IntersectionLattice(HyperplaneArrangement([[1,0],[0,1],[1,1]]));|
  <Geometric lattice: 3 atoms, rank 2>
  !gapprompt@gap>| !gapinput@LIsBoolean(L);|
  false
\end{Verbatim}
 }

 

\subsection{\textcolor{Chapter }{LIsGeneric (for IsGeomLattice)}}
\logpage{[ 2, 5, 3 ]}\nobreak
\hyperdef{L}{X7ACA5D857F2A0913}{}
{\noindent\textcolor{FuncColor}{$\triangleright$\enspace\texttt{LIsGeneric({\mdseries\slshape L})\index{LIsGeneric@\texttt{LIsGeneric}!for IsGeomLattice}
\label{LIsGeneric:for IsGeomLattice}
}\hfill{\scriptsize (property)}}\\
\textbf{\indent Returns:\ }
true or false. 



 Determines, if L is generic, i.e. if L is irreducible and all proper intervals
are boolean. 
\begin{Verbatim}[commandchars=!@|,fontsize=\small,frame=single,label=Example]
  !gapprompt@gap>| !gapinput@A := HyperplaneArrangement([[1,0,0],[0,1,0],[0,0,1],[1,1,1]]);|
  <HyperplaneArrangement: 4 hyperplanes in 3-space>
  !gapprompt@gap>| !gapinput@L:=IntersectionLattice(A);|
  <Geometric lattice: 4 atoms, rank 3>
  !gapprompt@gap>| !gapinput@LIsGeneric(L);|
  true
  !gapprompt@gap>| !gapinput@A := HyperplaneArrangement([[1,0,0],[0,1,0],[0,0,1],[1,1,0],[1,1,1]]);|
  <HyperplaneArrangement: 5 hyperplanes in 3-space>
  !gapprompt@gap>| !gapinput@L:=IntersectionLattice(A);|
  <Geometric lattice: 5 atoms, rank 3>
  !gapprompt@gap>| !gapinput@LIsGeneric(L);|
  false
\end{Verbatim}
 }

 }

 
\section{\textcolor{Chapter }{Operations}}\label{Chapter_Hyperplane_arrangements_Section_Operations}
\logpage{[ 2, 6, 0 ]}
\hyperdef{L}{X7DE8E16C7C2D387B}{}
{
  

\subsection{\textcolor{Chapter }{IsLEquiv (for IsHyperplaneArrangement,IsHyperplaneArrangement)}}
\logpage{[ 2, 6, 1 ]}\nobreak
\hyperdef{L}{X8367DC267DCC0635}{}
{\noindent\textcolor{FuncColor}{$\triangleright$\enspace\texttt{IsLEquiv({\mdseries\slshape A, B})\index{IsLEquiv@\texttt{IsLEquiv}!for IsHyperplaneArrangement,IsHyperplaneArrangement}
\label{IsLEquiv:for IsHyperplaneArrangement,IsHyperplaneArrangement}
}\hfill{\scriptsize (operation)}}\\
\textbf{\indent Returns:\ }
true or false 



 Determines if the arrangements \mbox{\texttt{\mdseries\slshape A}} and \mbox{\texttt{\mdseries\slshape B}} have isomorphic intersection lattices. 

 
\begin{Verbatim}[commandchars=!@|,fontsize=\small,frame=single,label=Example]
  !gapprompt@gap>| !gapinput@A:=Arr([[1,0,0],[0,1,0],[0,0,1],[1,1,0],[0,1,1],[1,1,3]]);|
  <HyperplaneArrangement: 6 hyperplanes in 3-space>
  !gapprompt@gap>| !gapinput@B:=AGpql(2,2,3);|
  <HyperplaneArrangement: 6 hyperplanes in 3-space>
  !gapprompt@gap>| !gapinput@IsLEquiv(A,B);|
  false
  !gapprompt@gap>| !gapinput@A:=Arr([[1,0,0],[0,1,0],[0,0,1],[1,1,0],[0,1,1],[1,1,1]]);|
  <HyperplaneArrangement: 6 hyperplanes in 3-space>
  !gapprompt@gap>| !gapinput@IsLEquiv(A,B);|
  true
\end{Verbatim}
 }

 

\subsection{\textcolor{Chapter }{Essentialization (for IsHyperplaneArrangement)}}
\logpage{[ 2, 6, 2 ]}\nobreak
\hyperdef{L}{X7BD35FA07804BB70}{}
{\noindent\textcolor{FuncColor}{$\triangleright$\enspace\texttt{Essentialization({\mdseries\slshape A})\index{Essentialization@\texttt{Essentialization}!for IsHyperplaneArrangement}
\label{Essentialization:for IsHyperplaneArrangement}
}\hfill{\scriptsize (operation)}}\\
\textbf{\indent Returns:\ }
A hyperplane arrangement. 



 Computes the \emph{essentialization} of the hyperplane arrangement \mbox{\texttt{\mdseries\slshape A}}. An arrangement is called essential if the intersection of all its
hyperplanes is the origin. If this is not the case, the function restricts the
arrangement to a complementary subspace so that the resulting arrangement
becomes essential. 
\begin{Verbatim}[commandchars=!@|,fontsize=\small,frame=single,label=Example]
  !gapprompt@gap>| !gapinput@A:=Arr([[1,0,0],[0,1,0],[1,1,0]]);|
  <HyperplaneArrangement: 3 hyperplanes in 3-space>
  !gapprompt@gap>| !gapinput@B:=Essentialization(A); Roots(B);|
  <HyperplaneArrangement: 3 hyperplanes in 2-space>
  [ [ 1, 0 ], [ 0, 1 ], [ 1, 1 ] ]
\end{Verbatim}
 }

 }

 
\section{\textcolor{Chapter }{Global methods}}\label{Chapter_Hyperplane_arrangements_Section_Global_methods}
\logpage{[ 2, 7, 0 ]}
\hyperdef{L}{X80841F8D80722081}{}
{
  

\subsection{\textcolor{Chapter }{HArrResHind}}
\logpage{[ 2, 7, 1 ]}\nobreak
\hyperdef{L}{X7CC338E67CE223FE}{}
{\noindent\textcolor{FuncColor}{$\triangleright$\enspace\texttt{HArrResHind({\mdseries\slshape A, i})\index{HArrResHind@\texttt{HArrResHind}}
\label{HArrResHind}
}\hfill{\scriptsize (function)}}\\
\textbf{\indent Returns:\ }
A hyperplane arrangement. 



 Computes the \emph{restriction} of the arrangement \mbox{\texttt{\mdseries\slshape A}} to the \mbox{\texttt{\mdseries\slshape i}}\texttt{\symbol{45}}th hyperplane of \texttt{Roots(A)}. }

 

\subsection{\textcolor{Chapter }{HArrResX}}
\logpage{[ 2, 7, 2 ]}\nobreak
\hyperdef{L}{X784AA86286AA9DD8}{}
{\noindent\textcolor{FuncColor}{$\triangleright$\enspace\texttt{HArrResX({\mdseries\slshape A, S})\index{HArrResX@\texttt{HArrResX}}
\label{HArrResX}
}\hfill{\scriptsize (function)}}\\
\textbf{\indent Returns:\ }
A hyperplane arrangement. 



 Computes the \emph{restriction} of the arrangement \mbox{\texttt{\mdseries\slshape A}} to the subspace orthogonal to the vectors in \texttt{S}. }

 

\subsection{\textcolor{Chapter }{DrawLatex3Arr}}
\logpage{[ 2, 7, 3 ]}\nobreak
\hyperdef{L}{X82C8EAF88340865B}{}
{\noindent\textcolor{FuncColor}{$\triangleright$\enspace\texttt{DrawLatex3Arr({\mdseries\slshape A[, ps[, ip[, Hind[, disthv[, MarkHs]]]]]})\index{DrawLatex3Arr@\texttt{DrawLatex3Arr}}
\label{DrawLatex3Arr}
}\hfill{\scriptsize (function)}}\\
\textbf{\indent Returns:\ }
A string. 



 Generates LaTeX tikz\texttt{\symbol{45}}code for a nice projective picture of
the real 3\texttt{\symbol{45}}arrangement. If $x,y,z$ are the coordinates, by default, this is the 2\texttt{\symbol{45}}dim affine
arrangement obtained by intersecting \mbox{\texttt{\mdseries\slshape A}} with the plane $z=1$. 

 The optional parameters are: 
\begin{itemize}
\item  \mbox{\texttt{\mdseries\slshape ps}} a scaling factor, 
\item  \mbox{\texttt{\mdseries\slshape ip}} intersection points are drawn, 
\item  \mbox{\texttt{\mdseries\slshape Hind}} labels for the hyperplanes are added, 
\item  \mbox{\texttt{\mdseries\slshape disthv}} a vector giving the normal of the plane whith which to intersect \mbox{\texttt{\mdseries\slshape A}}, 
\item  \mbox{\texttt{\mdseries\slshape MarkHs}} a list of indices of hyperplanes of \mbox{\texttt{\mdseries\slshape A}}, these are drawn in another color. 
\end{itemize}
 

 The example below will look as follows (only in pdf). 

 
\begin{Verbatim}[commandchars=!@|,fontsize=\small,frame=single,label=Example]
  !gapprompt@gap>| !gapinput@A:=AGpql(2,2,3);|
  <HyperplaneArrangement: 6 hyperplanes in 3-space>
  !gapprompt@gap>| !gapinput@Print(DrawLatex3Arr(A));|
  \begin{tikzpicture}[scale=1.0]
  \draw (-2.827,2.827) -- (2.827,-2.827);  % H_1 
  \draw (2.827,2.827) -- (-2.827,-2.827);  % H_2 
  \draw (-1.,3.873) -- (-1.,-3.873);  % H_3 
  \draw (1.,3.873) -- (1.,-3.873);  % H_4 
  \draw (3.873,-1.) -- (-3.873,-1.);  % H_5 
  \draw (3.873,1.) -- (-3.873,1.);  % H_6 
  \end{tikzpicture}
  !gapprompt@gap>| !gapinput@Print(DrawLatex3Arr(A,1/2,true,true,[1,1,1],[1,2]));|
  \begin{tikzpicture}[scale=1.0]
  \draw[color=red] (-3.559,1.827) -- (1.827,-3.559);  % H_1 
  \node at (1.975,-3.707) {\small $1$}; 
  \draw[color=red] (2.827,2.827) -- (-2.827,-2.827);  % H_2 
  \node at (-2.970,-2.970) {\small $2$}; 
  \draw (3.361,2.169) -- (-3.995,0.197);  % H_3 
  \node at (-4.198,0.143) {\small $3$}; 
  \draw (-1.034,3.863) -- (1.034,-3.863);  % H_4 
  \node at (1.087,-4.057) {\small $4$}; 
  \draw (2.169,3.361) -- (0.197,-3.995);  % H_5 
  \node at (0.143,-4.198) {\small $5$}; 
  \draw (-3.863,1.034) -- (3.863,-1.034);  % H_6 
  \node at (4.057,-1.087) {\small $6$}; 
  
  \fill[red] (-0.865,-0.865) circle[radius=2pt];  % P[ 1, 2 ] 
  \fill[red] (-2.366,0.634) circle[radius=2pt];  % P[ 1, 3, 6 ] 
  \fill[red] (1.732,1.732) circle[radius=2pt];  % P[ 2, 3, 5 ] 
  \fill[red] (0.634,-2.366) circle[radius=2pt];  % P[ 1, 4, 5 ] 
  \fill[red] (0.0,0.0) circle[radius=2pt];  % P[ 2, 4, 6 ] 
  \fill[red] (-0.317,1.183) circle[radius=2pt];  % P[ 3, 4 ] 
  \fill[red] (1.183,-0.317) circle[radius=2pt];  % P[ 5, 6 ] 
  
  \end{tikzpicture}
  
\end{Verbatim}
 

 
 \includegraphics{./other/DrawLatex3Arr_Example.pdf}
 

 }

 }

 }

   
\chapter{\textcolor{Chapter }{Free arrangements}}\label{Chapter_Free_arrangements}
\logpage{[ 3, 0, 0 ]}
\hyperdef{L}{X84CDDA0382B013CF}{}
{
  

 The following describes functions to analyse freeness properties of
arrangements such as \emph{inductive freeness} due to H. Terao \cite{Terao1980_FreeI} and \emph{divisionell freeness} due to T. Abe \cite{Abe16_DivFree}. 

 
\section{\textcolor{Chapter }{Freeness properties}}\label{Chapter_Free_arrangements_Section_Freeness_properties}
\logpage{[ 3, 1, 0 ]}
\hyperdef{L}{X82FF317F7EFD52AF}{}
{
  

\subsection{\textcolor{Chapter }{IsInductivelyFree (for IsHyperplaneArrangement)}}
\logpage{[ 3, 1, 1 ]}\nobreak
\hyperdef{L}{X80F4D6247CAA1BDA}{}
{\noindent\textcolor{FuncColor}{$\triangleright$\enspace\texttt{IsInductivelyFree({\mdseries\slshape A})\index{IsInductivelyFree@\texttt{IsInductivelyFree}!for IsHyperplaneArrangement}
\label{IsInductivelyFree:for IsHyperplaneArrangement}
}\hfill{\scriptsize (property)}}\\
\textbf{\indent Returns:\ }
true or false. 



 Determines whether the arrangement \mbox{\texttt{\mdseries\slshape A}} is \emph{inductively free}. 
\begin{Verbatim}[commandchars=!@|,fontsize=\small,frame=single,label=Example]
  !gapprompt@gap>| !gapinput@A:=HyperplaneArrangement([[1,0,0],[0,1,0],[0,0,1],[1,1,0],[1,1,1]]);|
  <HyperplaneArrangement: 5 hyperplanes in 3-space>
  !gapprompt@gap>| !gapinput@IsInductivelyFree(A);|
  true
\end{Verbatim}
 }

 

\subsection{\textcolor{Chapter }{IsDivisionallyFree (for IsHyperplaneArrangement)}}
\logpage{[ 3, 1, 2 ]}\nobreak
\hyperdef{L}{X85DAC6FC78DDED6D}{}
{\noindent\textcolor{FuncColor}{$\triangleright$\enspace\texttt{IsDivisionallyFree({\mdseries\slshape A})\index{IsDivisionallyFree@\texttt{IsDivisionallyFree}!for IsHyperplaneArrangement}
\label{IsDivisionallyFree:for IsHyperplaneArrangement}
}\hfill{\scriptsize (property)}}\\
\textbf{\indent Returns:\ }
true or false. 



 Determines whether the arrangement \mbox{\texttt{\mdseries\slshape A}} is \emph{divisionally free}. 
\begin{Verbatim}[commandchars=!@|,fontsize=\small,frame=single,label=Example]
  !gapprompt@gap>| !gapinput@A:=HyperplaneArrangement([[1,0,0],[0,1,0],[0,0,1],[1,1,0],[1,1,1]]);|
  <HyperplaneArrangement: 5 hyperplanes in 3-space>
  !gapprompt@gap>| !gapinput@IsDivisionallyFree(A);|
  true
  !gapprompt@gap>| !gapinput@A:=AGpql(3,3,3);|
  <HyperplaneArrangement: 9 hyperplanes in 3-space>
  !gapprompt@gap>| !gapinput@IsDivisionallyFree(A);|
  false
\end{Verbatim}
 }

 }

 }

   
\chapter{\textcolor{Chapter }{Milnor fibers}}\label{Chapter_Milnor_fibers}
\logpage{[ 4, 0, 0 ]}
\hyperdef{L}{X80236F007DAFE4C2}{}
{
  

 The Milnor fiber is the typical fiber $f^{-1}(1)$ of the Milnor fibration 
\[f:\mathbb{C}^\ell \setminus \cup_{H \in \mathcal{A}}H \to \mathbb{C}^\times,
z\mapsto \prod_{H \in \mathcal{A}}\alpha_H(z)\]
 of the hyperplane arrangement $\mathcal{A}$, where $\alpha_H$ are defining linear forms. 

 The following describes functions to construct a regular cell complex having
the homotopy type of the Milnor fiber from \cite{MueYos25_MilnorOM} and thus enables computation of homotopy invaraints using e.g. the HAP package \cite{GAP_HAP}. 
\section{\textcolor{Chapter }{Complexes}}\label{Chapter_Milnor_fibers_Section_Complexes}
\logpage{[ 4, 1, 0 ]}
\hyperdef{L}{X79C388D385DB7CD1}{}
{
  

\subsection{\textcolor{Chapter }{MilnorFiberComplex (for IsOrientedMatroid)}}
\logpage{[ 4, 1, 1 ]}\nobreak
\hyperdef{L}{X794B90A67B153E6F}{}
{\noindent\textcolor{FuncColor}{$\triangleright$\enspace\texttt{MilnorFiberComplex({\mdseries\slshape OM})\index{MilnorFiberComplex@\texttt{MilnorFiberComplex}!for IsOrientedMatroid}
\label{MilnorFiberComplex:for IsOrientedMatroid}
}\hfill{\scriptsize (attribute)}}\\
\textbf{\indent Returns:\ }
FacePoset 



 Computes the face poset of the regular cell complex of the combinatorial
Milnor fiber of the oriented matroid \mbox{\texttt{\mdseries\slshape OM}}. 
\begin{Verbatim}[commandchars=!@|,fontsize=\small,frame=single,label=Example]
  !gapprompt@gap>| !gapinput@O:=OrientedMatroid([[1,0],[0,1],[1,1]]);|
  <OrientedMatroid: 3 elements, rank 2>
  !gapprompt@gap>| !gapinput@MCpx:=MilnorFiberComplex(O);|
  <FacePoset of dimension 1 with f-vector [ 3, 6 ]>
\end{Verbatim}
 }

 

\subsection{\textcolor{Chapter }{MilnorFiberComplex (for IsHyperplaneArrangement)}}
\logpage{[ 4, 1, 2 ]}\nobreak
\hyperdef{L}{X7D30E85086B93DB0}{}
{\noindent\textcolor{FuncColor}{$\triangleright$\enspace\texttt{MilnorFiberComplex({\mdseries\slshape A})\index{MilnorFiberComplex@\texttt{MilnorFiberComplex}!for IsHyperplaneArrangement}
\label{MilnorFiberComplex:for IsHyperplaneArrangement}
}\hfill{\scriptsize (attribute)}}\\
\textbf{\indent Returns:\ }
FacePoset 



 Computes the face poset of the regular cell complex having the homotopy type
of Milnor fiber of the arrangement \mbox{\texttt{\mdseries\slshape A}}. 
\begin{Verbatim}[commandchars=!@|,fontsize=\small,frame=single,label=Example]
  !gapprompt@gap>| !gapinput@A:=AGpql(2,2,3);|
  <HyperplaneArrangement: 6 hyperplanes in 3-space>
  !gapprompt@gap>| !gapinput@MCpx:=MilnorFiberComplex(A);|
  <FacePoset of dimension 2 with f-vector [ 12, 60, 60 ]>
\end{Verbatim}
 }

 

\subsection{\textcolor{Chapter }{FPtoCWCpx}}
\logpage{[ 4, 1, 3 ]}\nobreak
\hyperdef{L}{X7BBA6F967E493B5D}{}
{\noindent\textcolor{FuncColor}{$\triangleright$\enspace\texttt{FPtoCWCpx({\mdseries\slshape FP})\index{FPtoCWCpx@\texttt{FPtoCWCpx}}
\label{FPtoCWCpx}
}\hfill{\scriptsize (function)}}\\
\textbf{\indent Returns:\ }
HAP CW Complex 



 Converts a \textsc{HypArr} \texttt{FacePoset} into a \textsc{HAP} \texttt{RegularCWComplex} for further computations of topological invariants. 
\begin{Verbatim}[commandchars=!@|,fontsize=\small,frame=single,label=Example]
  !gapprompt@gap>| !gapinput@A:=AGpql(2,2,3);|
  <HyperplaneArrangement: 6 hyperplanes in 3-space>
  !gapprompt@gap>| !gapinput@MCpx:=MilnorFiberComplex(A);|
  <FacePoset of dimension 2 with f-vector [ 12, 60, 60 ]>
  !gapprompt@gap>| !gapinput@MCW := FPtoCWCpx(MCpx);|
  Regular CW-complex of dimension 2
  !gapprompt@gap>| !gapinput@Homology(MCW,0);|
  [ 0 ]
  !gapprompt@gap>| !gapinput@Homology(MCW,1);|
  [ 0, 0, 0, 0, 0, 0, 0 ]
\end{Verbatim}
 }

 }

 }

   
\chapter{\textcolor{Chapter }{Further properties of arrangements, geometric lattices, and oriented matroids}}\label{Chapter_Further_properties_of_arrangements_geometric_lattices_and_oriented_matroids}
\logpage{[ 5, 0, 0 ]}
\hyperdef{L}{X7E898D0F86EF9F1E}{}
{
  

 The following describes functions to analyse further properties of
arrangements such as \emph{formality}, \emph{supersolvability}, \emph{simpliciality}, \emph{factoredness}, \emph{inductive factoredness} ... 

 
\section{\textcolor{Chapter }{Supersolvable arrangements}}\label{Chapter_Further_properties_of_arrangements_geometric_lattices_and_oriented_matroids_Section_Supersolvable_arrangements}
\logpage{[ 5, 1, 0 ]}
\hyperdef{L}{X7B09D56B7A076E6B}{}
{
  

\subsection{\textcolor{Chapter }{LIsModularPair (for IsGeomLattice, IsList, IsList)}}
\logpage{[ 5, 1, 1 ]}\nobreak
\hyperdef{L}{X78E6E11179F7ECB5}{}
{\noindent\textcolor{FuncColor}{$\triangleright$\enspace\texttt{LIsModularPair({\mdseries\slshape L, m1, m2})\index{LIsModularPair@\texttt{LIsModularPair}!for IsGeomLattice, IsList, IsList}
\label{LIsModularPair:for IsGeomLattice, IsList, IsList}
}\hfill{\scriptsize (operation)}}\\
\textbf{\indent Returns:\ }
true or false 



 Determines if \mbox{\texttt{\mdseries\slshape m1}} and \mbox{\texttt{\mdseries\slshape m2}} form a modular pair in L. 
\begin{Verbatim}[commandchars=!@|,fontsize=\small,frame=single,label=Example]
  !gapprompt@gap>| !gapinput@A:=AGpql(2,2,3);|
  <HyperplaneArrangement: 6 hyperplanes in 3-space>
  !gapprompt@gap>| !gapinput@L:=IntersectionLattice(A);|
  <Geometric lattice: 6 atoms, rank 3>
  !gapprompt@gap>| !gapinput@LGroundSet(L);|
  [ [ [ 1 ], [ 2 ], [ 3 ], [ 4 ], [ 5 ], [ 6 ] ], 
   [ [ 1, 2 ], [ 1, 3, 6 ], [ 2, 3, 5 ], [ 1, 4, 5 ], 
     [ 2, 4, 6 ], [ 3, 4 ], [ 5, 6 ] ], 
    [ [ 1, 2, 3, 4, 5, 6 ] ] ]
  !gapprompt@gap>| !gapinput@m1:=[1,2];; m2:= [5,6];; LIsModularPair(L,m1,m2);|
  false
  !gapprompt@gap>| !gapinput@m1:=[1,2];; m2:= [2,4,6];; LIsModularPair(L,m1,m2);|
  true
\end{Verbatim}
 }

 

\subsection{\textcolor{Chapter }{LIsModularFlat (for IsGeomLattice, IsList)}}
\logpage{[ 5, 1, 2 ]}\nobreak
\hyperdef{L}{X7FECD913876C637C}{}
{\noindent\textcolor{FuncColor}{$\triangleright$\enspace\texttt{LIsModularFlat({\mdseries\slshape L, m})\index{LIsModularFlat@\texttt{LIsModularFlat}!for IsGeomLattice, IsList}
\label{LIsModularFlat:for IsGeomLattice, IsList}
}\hfill{\scriptsize (operation)}}\\
\textbf{\indent Returns:\ }
true or false 



 Determines if \mbox{\texttt{\mdseries\slshape m}} is a modular flat in L. 
\begin{Verbatim}[commandchars=!@|,fontsize=\small,frame=single,label=Example]
  !gapprompt@gap>| !gapinput@A:=AGpql(2,2,3); L:=IntersectionLattice(A); LGroundSet(L);|
  <HyperplaneArrangement: 6 hyperplanes in 3-space>
  <Geometric lattice: 6 atoms, rank 3>
  [ [ [ 1 ], [ 2 ], [ 3 ], [ 4 ], [ 5 ], [ 6 ] ], 
   [ [ 1, 2 ], [ 1, 3, 6 ], [ 2, 3, 5 ], [ 1, 4, 5 ], 
     [ 2, 4, 6 ], [ 3, 4 ], [ 5, 6 ] ], 
    [ [ 1, 2, 3, 4, 5, 6 ] ] ]
  !gapprompt@gap>| !gapinput@m1:=[1,2];; m2:= [2,4,6];; LIsModularFlat(L,m1); LIsModularFlat(L,m2);|
  false
  true
\end{Verbatim}
 }

 

\subsection{\textcolor{Chapter }{LModularFlatsRk (for IsGeomLattice, IsInt)}}
\logpage{[ 5, 1, 3 ]}\nobreak
\hyperdef{L}{X8271B52F816053DB}{}
{\noindent\textcolor{FuncColor}{$\triangleright$\enspace\texttt{LModularFlatsRk({\mdseries\slshape L, k})\index{LModularFlatsRk@\texttt{LModularFlatsRk}!for IsGeomLattice, IsInt}
\label{LModularFlatsRk:for IsGeomLattice, IsInt}
}\hfill{\scriptsize (operation)}}\\
\textbf{\indent Returns:\ }
list 



 Determines the modular flats in \mbox{\texttt{\mdseries\slshape L}} of rank \mbox{\texttt{\mdseries\slshape k}}. 
\begin{Verbatim}[commandchars=!@|,fontsize=\small,frame=single,label=Example]
  !gapprompt@gap>| !gapinput@L:=IntersectionLattice(AGpql(2,1,4));|
  <Geometric lattice: 16 atoms, rank 4>
  !gapprompt@gap>| !gapinput@LModularFlatsRk(L,3);|
  [ [ 1, 2, 3, 5, 6, 7, 8, 11, 12 ], 
    [ 1, 2, 4, 5, 6, 9, 10, 13, 14 ], 
    [ 1, 3, 4, 7, 8, 9, 10, 15, 16 ], 
    [ 2, 3, 4, 11, 12, 13, 14, 15, 16 ] ]
\end{Verbatim}
 }

 

\subsection{\textcolor{Chapter }{LIsSupersolvable (for IsGeomLattice)}}
\logpage{[ 5, 1, 4 ]}\nobreak
\hyperdef{L}{X7D95D0F37E6FBFD6}{}
{\noindent\textcolor{FuncColor}{$\triangleright$\enspace\texttt{LIsSupersolvable({\mdseries\slshape L})\index{LIsSupersolvable@\texttt{LIsSupersolvable}!for IsGeomLattice}
\label{LIsSupersolvable:for IsGeomLattice}
}\hfill{\scriptsize (property)}}\\
\textbf{\indent Returns:\ }
true or false. 



 Determines whether the geometric lattice \mbox{\texttt{\mdseries\slshape L}} is \emph{supersolvable}. 
\begin{Verbatim}[commandchars=!@|,fontsize=\small,frame=single,label=Example]
  !gapprompt@gap>| !gapinput@L:=IntersectionLattice(AGpql(2,1,4));|
  <Geometric lattice: 16 atoms, rank 4>
  !gapprompt@gap>| !gapinput@LIsSupersolvable(L);|
  true
  !gapprompt@gap>| !gapinput@L:=IntersectionLattice(AGpql(2,2,4));|
  <Geometric lattice: 12 atoms, rank 4>
  !gapprompt@gap>| !gapinput@LIsSupersolvable(L);|
  false
\end{Verbatim}
 }

 

\subsection{\textcolor{Chapter }{HArrIsSupersolvable (for IsHyperplaneArrangement)}}
\logpage{[ 5, 1, 5 ]}\nobreak
\hyperdef{L}{X7DC90A0286D5664A}{}
{\noindent\textcolor{FuncColor}{$\triangleright$\enspace\texttt{HArrIsSupersolvable({\mdseries\slshape A})\index{HArrIsSupersolvable@\texttt{HArrIsSupersolvable}!for IsHyperplaneArrangement}
\label{HArrIsSupersolvable:for IsHyperplaneArrangement}
}\hfill{\scriptsize (property)}}\\
\textbf{\indent Returns:\ }
true or false. 



 Determines whether the arrangement \mbox{\texttt{\mdseries\slshape A}} is \emph{supersolvable}. 
\begin{Verbatim}[commandchars=!@|,fontsize=\small,frame=single,label=Example]
  !gapprompt@gap>| !gapinput@A:=AGpql(2,1,4);|
  <HyperplaneArrangement: 16 hyperplanes in 4-space>
  !gapprompt@gap>| !gapinput@HArrIsSupersolvable(A);|
  true
  !gapprompt@gap>| !gapinput@A:=AGpql(2,2,4); |
  <HyperplaneArrangement: 12 hyperplanes in 4-space>
  !gapprompt@gap>| !gapinput@HArrIsSupersolvable(A);|
  false
\end{Verbatim}
 }

 

\subsection{\textcolor{Chapter }{OMIsSupersolvable (for IsOrientedMatroid)}}
\logpage{[ 5, 1, 6 ]}\nobreak
\hyperdef{L}{X78B7232883025D41}{}
{\noindent\textcolor{FuncColor}{$\triangleright$\enspace\texttt{OMIsSupersolvable({\mdseries\slshape OM})\index{OMIsSupersolvable@\texttt{OMIsSupersolvable}!for IsOrientedMatroid}
\label{OMIsSupersolvable:for IsOrientedMatroid}
}\hfill{\scriptsize (property)}}\\
\textbf{\indent Returns:\ }
true or false. 



 Determines whether the oriented matroid \mbox{\texttt{\mdseries\slshape OM}} is \emph{supersolvable}. 
\begin{Verbatim}[commandchars=!@|,fontsize=\small,frame=single,label=Example]
  !gapprompt@gap>| !gapinput@O:=OrientedMatroid(3,4,[1,1,1,1]);|
  <OrientedMatroid: 4 elements, rank 3>
  !gapprompt@gap>| !gapinput@OMIsSupersolvable(O);|
  false
  !gapprompt@gap>| !gapinput@C:=List(Combinations(Roots(AGpql(2,2,3)),3),x->pos(Determinant(x)));|
  [ 1, 1, 1, 1, 0, -1, -1, -1, -1, 0, 1, 0, 1, -1, 0, 1, 1, 1, -1, -1 ]
  !gapprompt@gap>| !gapinput@O:=OrientedMatroid(3,6,C);|
  <OrientedMatroid: 6 elements, rank 3>
  !gapprompt@gap>| !gapinput@OMIsSupersolvable(O);|
  true
\end{Verbatim}
 }

 }

 
\section{\textcolor{Chapter }{Formal arrangements}}\label{Chapter_Further_properties_of_arrangements_geometric_lattices_and_oriented_matroids_Section_Formal_arrangements}
\logpage{[ 5, 2, 0 ]}
\hyperdef{L}{X79BAC87F7C1CA4A1}{}
{
  

\subsection{\textcolor{Chapter }{IsFormal (for IsHyperplaneArrangement)}}
\logpage{[ 5, 2, 1 ]}\nobreak
\hyperdef{L}{X808977C7824120DF}{}
{\noindent\textcolor{FuncColor}{$\triangleright$\enspace\texttt{IsFormal({\mdseries\slshape A})\index{IsFormal@\texttt{IsFormal}!for IsHyperplaneArrangement}
\label{IsFormal:for IsHyperplaneArrangement}
}\hfill{\scriptsize (property)}}\\
\textbf{\indent Returns:\ }
true or false. 



 Determines whether the arrangement is \emph{formal} in the sense of Falk\texttt{\symbol{45}}Randell \cite{FalkRandell87_HomotopyArr}. 
\begin{Verbatim}[commandchars=!@|,fontsize=\small,frame=single,label=Example]
  !gapprompt@gap>| !gapinput@A1 := Arr([|
  !gapprompt@>| !gapinput@    [1,0,0],|
  !gapprompt@>| !gapinput@    [0,1,0],|
  !gapprompt@>| !gapinput@    [0,0,1],|
  !gapprompt@>| !gapinput@    [1,1,1],|
  !gapprompt@>| !gapinput@    [2,1,1],|
  !gapprompt@>| !gapinput@    [2,3,1],|
  !gapprompt@>| !gapinput@    [2,3,4],|
  !gapprompt@>| !gapinput@    [3,0,5],|
  !gapprompt@>| !gapinput@    [3,4,5]|
  !gapprompt@>| !gapinput@]);;|
  !gapprompt@gap>| !gapinput@|
  !gapprompt@gap>| !gapinput@A2 := Arr([|
  !gapprompt@>| !gapinput@    [1,0,0],|
  !gapprompt@>| !gapinput@    [0,1,0],|
  !gapprompt@>| !gapinput@    [0,0,1],|
  !gapprompt@>| !gapinput@    [1,1,1],|
  !gapprompt@>| !gapinput@    [2,1,1],|
  !gapprompt@>| !gapinput@    [2,3,1],|
  !gapprompt@>| !gapinput@    [2,3,4],|
  !gapprompt@>| !gapinput@    [1,0,3],|
  !gapprompt@>| !gapinput@    [1,2,3]|
  !gapprompt@>| !gapinput@]);;|
  !gapprompt@gap>| !gapinput@IsLEquiv(A1,A2);|
  true
  !gapprompt@gap>| !gapinput@IsFormal(A1);|
  true
  !gapprompt@gap>| !gapinput@IsFormal(A2);|
  false
\end{Verbatim}
 }

 }

 
\section{\textcolor{Chapter }{Simplicial arrangements}}\label{Chapter_Further_properties_of_arrangements_geometric_lattices_and_oriented_matroids_Section_Simplicial_arrangements}
\logpage{[ 5, 3, 0 ]}
\hyperdef{L}{X8070A6CF85A00A70}{}
{
  

\subsection{\textcolor{Chapter }{HArrIsSimplicial (for IsHyperplaneArrangement)}}
\logpage{[ 5, 3, 1 ]}\nobreak
\hyperdef{L}{X7D69068B7DA44D16}{}
{\noindent\textcolor{FuncColor}{$\triangleright$\enspace\texttt{HArrIsSimplicial({\mdseries\slshape A})\index{HArrIsSimplicial@\texttt{HArrIsSimplicial}!for IsHyperplaneArrangement}
\label{HArrIsSimplicial:for IsHyperplaneArrangement}
}\hfill{\scriptsize (property)}}\\
\textbf{\indent Returns:\ }
true or false. 



 Determines whether the real arrangement \mbox{\texttt{\mdseries\slshape A}} is \emph{simplicial}, i.e. if all chambers are simplicial. 
\begin{Verbatim}[commandchars=!@|,fontsize=\small,frame=single,label=Example]
  !gapprompt@gap>| !gapinput@A:=HyperplaneArrangement([[1,0,0],[0,1,0],[0,0,1],[1,1,0],[0,1,1]]); HArrIsSimplicial(A);|
  <HyperplaneArrangement: 5 hyperplanes in 3-space>
  false
  !gapprompt@gap>| !gapinput@A:=HyperplaneArrangement([[1,0,0],[0,1,0],[0,0,1],[1,1,0],[0,1,1],[1,1,1]]); HArrIsSimplicial(A);|
  <HyperplaneArrangement: 6 hyperplanes in 3-space>
  true
\end{Verbatim}
 }

 }

 
\section{\textcolor{Chapter }{Falk's weight test}}\label{Chapter_Further_properties_of_arrangements_geometric_lattices_and_oriented_matroids_Section_Falks_weight_test}
\logpage{[ 5, 4, 0 ]}
\hyperdef{L}{X79A9BA6287B1A981}{}
{
  The following still needs to be tested further... 

\subsection{\textcolor{Chapter }{OMBoundedCpx (for IsOrientedMatroid, IsInt)}}
\logpage{[ 5, 4, 1 ]}\nobreak
\hyperdef{L}{X7F2F674F8273D120}{}
{\noindent\textcolor{FuncColor}{$\triangleright$\enspace\texttt{OMBoundedCpx({\mdseries\slshape OM, g})\index{OMBoundedCpx@\texttt{OMBoundedCpx}!for IsOrientedMatroid, IsInt}
\label{OMBoundedCpx:for IsOrientedMatroid, IsInt}
}\hfill{\scriptsize (operation)}}\\
\textbf{\indent Returns:\ }
list of list of sign vectors. 



 Constructs the complex of bounded cells of the affine part of the oriented
matroid \mbox{\texttt{\mdseries\slshape OM}} with respect to the element \mbox{\texttt{\mdseries\slshape g}}. 
\begin{Verbatim}[commandchars=!@|,fontsize=\small,frame=single,label=Example]
  
\end{Verbatim}
 }

 

\subsection{\textcolor{Chapter }{OMSupportsFalkWeights (for IsOrientedMatroid, IsInt)}}
\logpage{[ 5, 4, 2 ]}\nobreak
\hyperdef{L}{X87CE42827934F11F}{}
{\noindent\textcolor{FuncColor}{$\triangleright$\enspace\texttt{OMSupportsFalkWeights({\mdseries\slshape OM, g})\index{OMSupportsFalkWeights@\texttt{OMSupportsFalkWeights}!for IsOrientedMatroid, IsInt}
\label{OMSupportsFalkWeights:for IsOrientedMatroid, IsInt}
}\hfill{\scriptsize (operation)}}\\
\textbf{\indent Returns:\ }
list or fail 



 Only for oriented matroids of rank 3. Determines, whether \mbox{\texttt{\mdseries\slshape OM}} supports a system of weights as described by M. Falk in \cite{Falk1995_Kpi1Arrs}. This implies, that the Salvetti complex (see \texttt{SalvettiComplex} (\ref{SalvettiComplex:for IsOrientedMatroid})) is apsherical. The functionality of the \textsc{cddinterface} package is used, to determine solotions of a system of linear inequalities. 

 If the algorithm determines a suitable weight system for \mbox{\texttt{\mdseries\slshape OM}}, each list element consists of a weight, and a corner, i.e. a pair of a
2\texttt{\symbol{45}}cell and an adjacent vertex in the bounded complex, given
as sign vectors. 
\begin{Verbatim}[commandchars=!@|,fontsize=\small,frame=single,label=Example]
  !gapprompt@gap>| !gapinput@O:=OrientedMatroid(|
  !gapprompt@>| !gapinput@      [[1,1,0],[1,-1,0],[1,0,1],[1,0,-1],[0,1,1],[0,1,-1],[0,0,1]]|
  !gapprompt@>| !gapinput@);;|
  <OrientedMatroid: 7 elements, rank 3>
  !gapprompt@gap>| !gapinput@OMIsSupersolvable(O);|
  false
  !gapprompt@gap>| !gapinput@OMSupportsFalkWeights(O,3);|
  fail
  !gapprompt@gap>| !gapinput@OMSupportsFalkWeights(O,1);|
  [ [ 1/2, [ [ 1, 1, 1, 1, 1, 1, 1 ], [ 1, 0, 1, 0, 1, 0, 1 ] ] ], 
    [ 0, [ [ 1, 1, 1, 1, 1, 1, 1 ], [ 1, 1, 1, 1, 0, 0, 0 ] ] ], 
    [ 1/2, [ [ 1, 1, 1, 1, 1, 1, 1 ], [ 1, 0, 1, 1, 1, 1, 0 ] ] ], 
    [ 0, [ [ 1, 1, 1, 1, 1, 1, -1 ], [ 1, 0, 0, 1, 0, 1, -1 ] ] ], 
    [ 1/2, [ [ 1, 1, 1, 1, 1, 1, -1 ], [ 1, 1, 1, 1, 0, 0, 0 ] ] ], 
    [ 1/2, [ [ 1, 1, 1, 1, 1, 1, -1 ], [ 1, 0, 1, 1, 1, 1, 0 ] ] ], 
    [ 0, [ [ 1, -1, 1, 1, 1, 1, 1 ], [ 1, 0, 1, 0, 1, 0, 1 ] ] ], 
    [ 1/2, [ [ 1, -1, 1, 1, 1, 1, 1 ], [ 1, -1, 0, 0, 1, 1, 0 ] ] ], 
    [ 1/2, [ [ 1, -1, 1, 1, 1, 1, 1 ], [ 1, 0, 1, 1, 1, 1, 0 ] ] ], 
    [ 1/2, [ [ 1, -1, 1, 1, 1, 1, -1 ], [ 1, 0, 0, 1, 0, 1, -1 ] ] ], 
    [ 0, [ [ 1, -1, 1, 1, 1, 1, -1 ], [ 1, -1, 0, 0, 1, 1, 0 ] ] ], 
    [ 1/2, [ [ 1, -1, 1, 1, 1, 1, -1 ], [ 1, 0, 1, 1, 1, 1, 0 ] ] ] ] 
\end{Verbatim}
 }

 }

 }

   
\chapter{\textcolor{Chapter }{Oriented Matroids}}\label{Chapter_Oriented_Matroids}
\logpage{[ 6, 0, 0 ]}
\hyperdef{L}{X8388FBA579602CDB}{}
{
  

 Oriented matroids are abstractions of real hyperplane arrangements. The
following describes functions to handle oriented matroids and compute
associated cell complexes and invariants. See \cite{BLSWZ1999_OrientedMatroids}. 
\section{\textcolor{Chapter }{Construction of oriented matroids}}\label{Chapter_Oriented_Matroids_Section_Construction_of_oriented_matroids}
\logpage{[ 6, 1, 0 ]}
\hyperdef{L}{X8053D2C379B5CD9C}{}
{
  

\subsection{\textcolor{Chapter }{OrientedMatroid (for IsList)}}
\logpage{[ 6, 1, 1 ]}\nobreak
\hyperdef{L}{X82EE52377B03CC80}{}
{\noindent\textcolor{FuncColor}{$\triangleright$\enspace\texttt{OrientedMatroid({\mdseries\slshape R, or, A, or[, CCs, c], or[, n, r, ChiroCoreList]})\index{OrientedMatroid@\texttt{OrientedMatroid}!for IsList}
\label{OrientedMatroid:for IsList}
}\hfill{\scriptsize (operation)}}\\
\textbf{\indent Returns:\ }
An oriented matroid OM. 



 Constructs am oriented matroid from 
\begin{itemize}
\item  a list \mbox{\texttt{\mdseries\slshape R}} of vectors representing defining linear forms of real hyperplanes, 
\item  a real hyperplane arrangement \mbox{\texttt{\mdseries\slshape A}}, 
\item  a list of lists \mbox{\texttt{\mdseries\slshape CVs}} of signed covectors and a string \texttt{c="cv"}, 
\item  intergers \mbox{\texttt{\mdseries\slshape n,r}} giving the size of the groundset and the rank and a list \mbox{\texttt{\mdseries\slshape ChiroCoreList}} $\subseteq \{0,1,-1\}^m (m = \binom{n}{r})$ presenting the chirotope 
\end{itemize}
 

 of the oriented matroid. 
\begin{Verbatim}[commandchars=!@|,fontsize=\small,frame=single,label=Example]
  !gapprompt@gap>| !gapinput@A:=AGpql(2,2,3); OM:=OrientedMatroid(Roots(A));|
  <OrientedMatroid: 6 elements, rank 3>
  !gapprompt@gap>| !gapinput@A:=AGpql(2,2,3); OM:=OrientedMatroid(A);|
  <OrientedMatroid: 6 elements, rank 3>
  !gapprompt@gap>| !gapinput@OrientedMatroid([[[1,1],[-1,1],[1,-1],[-1,-1]],[[0,1],[0,-1],[1,0],[-1,0]],[[0,0]]],"cv");|
  <OrientedMatroid: 2 elements, rank 2>
  !gapprompt@gap>| !gapinput@OrientedMatroid(2,3,[1,1,1]);|
  <OrientedMatroid: 3 elements, rank 2>
\end{Verbatim}
 }

 }

 
\section{\textcolor{Chapter }{Attributes}}\label{Chapter_Oriented_Matroids_Section_Attributes}
\logpage{[ 6, 2, 0 ]}
\hyperdef{L}{X7C701DBF7BAE649A}{}
{
  

\subsection{\textcolor{Chapter }{OMRank (for IsOrientedMatroid)}}
\logpage{[ 6, 2, 1 ]}\nobreak
\hyperdef{L}{X7E30B85C7CE42908}{}
{\noindent\textcolor{FuncColor}{$\triangleright$\enspace\texttt{OMRank({\mdseries\slshape OM})\index{OMRank@\texttt{OMRank}!for IsOrientedMatroid}
\label{OMRank:for IsOrientedMatroid}
}\hfill{\scriptsize (attribute)}}\\
\textbf{\indent Returns:\ }
An integer 



 Returns the rank of the oriented matroid \mbox{\texttt{\mdseries\slshape OM}}. 
\begin{Verbatim}[commandchars=!@|,fontsize=\small,frame=single,label=Example]
  !gapprompt@gap>| !gapinput@A := AGpql(2,2,3); OM := OrientedMatroid(A);|
  <OrientedMatroid: 6 elements, rank 3>
  !gapprompt@gap>| !gapinput@OMRank(OM);|
  3
\end{Verbatim}
 }

 

\subsection{\textcolor{Chapter }{OMGroundSet (for IsOrientedMatroid)}}
\logpage{[ 6, 2, 2 ]}\nobreak
\hyperdef{L}{X78C0E4A0861CAD2F}{}
{\noindent\textcolor{FuncColor}{$\triangleright$\enspace\texttt{OMGroundSet({\mdseries\slshape OM})\index{OMGroundSet@\texttt{OMGroundSet}!for IsOrientedMatroid}
\label{OMGroundSet:for IsOrientedMatroid}
}\hfill{\scriptsize (attribute)}}\\
\textbf{\indent Returns:\ }
A list 



 Returns the ground set of the oriented matroid OM. 
\begin{Verbatim}[commandchars=!@|,fontsize=\small,frame=single,label=Example]
  !gapprompt@gap>| !gapinput@A := AGpql(2,2,3); OM := OrientedMatroid(A);|
  <OrientedMatroid: 6 elements, rank 3>
  !gapprompt@gap>| !gapinput@OMGroundSet(OM);|
  [ 1, 2, 3, 4, 5, 6 ]
\end{Verbatim}
 }

 

\subsection{\textcolor{Chapter }{OMLForms (for IsOrientedMatroid)}}
\logpage{[ 6, 2, 3 ]}\nobreak
\hyperdef{L}{X86AAB8F480E92ABB}{}
{\noindent\textcolor{FuncColor}{$\triangleright$\enspace\texttt{OMLForms({\mdseries\slshape OM})\index{OMLForms@\texttt{OMLForms}!for IsOrientedMatroid}
\label{OMLForms:for IsOrientedMatroid}
}\hfill{\scriptsize (attribute)}}\\
\textbf{\indent Returns:\ }
A list of linear forms or \texttt{fail} 



 Returns the list of linear forms associated with the oriented matroid OM. If
the internal component \texttt{OM!.lforms} is not bound, \texttt{fail} is returned. 
\begin{Verbatim}[commandchars=!@|,fontsize=\small,frame=single,label=Example]
  !gapprompt@gap>| !gapinput@A := AGpql(2,2,3); OM := OrientedMatroid(A);|
  <OrientedMatroid: 6 elements, rank 3>
  !gapprompt@gap>| !gapinput@OMLForms(O);|
  [ [ 1, 1, 0 ], [ 1, -1, 0 ], [ 1, 0, 1 ], 
    [ 1, 0, -1 ], [ 0, 1, 1 ], [ 0, 1, -1 ] ]
\end{Verbatim}
 }

 

\subsection{\textcolor{Chapter }{OMGeomLattice (for IsOrientedMatroid)}}
\logpage{[ 6, 2, 4 ]}\nobreak
\hyperdef{L}{X81CC016D818B3068}{}
{\noindent\textcolor{FuncColor}{$\triangleright$\enspace\texttt{OMGeomLattice({\mdseries\slshape OM})\index{OMGeomLattice@\texttt{OMGeomLattice}!for IsOrientedMatroid}
\label{OMGeomLattice:for IsOrientedMatroid}
}\hfill{\scriptsize (attribute)}}\\
\textbf{\indent Returns:\ }
An object of type \texttt{IsGeomLattice} or \texttt{fail} 



 Constructs the geometric lattice associated with the oriented matroid OM. If
OM is linear (as determined by \texttt{OMIsLinear(OM)}), the lattice is computed as the intersection lattice of the hyperplane
arrangement given by \texttt{OMLForms(OM)}, stored in the internal component \texttt{OM!.lattice}, and returned. 
\begin{Verbatim}[commandchars=!@|,fontsize=\small,frame=single,label=Example]
  !gapprompt@gap>| !gapinput@O := OrientedMatroid([[1,0,0],[0,1,0],[0,0,1],[1,1,1]]);|
  <OrientedMatroid: 4 elements, rank 3>
  !gapprompt@gap>| !gapinput@L:=OMGeomLattice(O); LGroundSet(L);|
  <Geometric lattice: 4 atoms, rank 3>
  [ [ [ 1 ], [ 2 ], [ 3 ], [ 4 ] ], 
    [ [ 1, 2 ], [ 1, 3 ], [ 2, 3 ], 
        [ 1, 4 ], [ 2, 4 ], [ 3, 4 ] ], 
    [ [ 1, 2, 3, 4 ] ] ]
  !gapprompt@gap>| !gapinput@O := OrientedMatroid(3,4,[1,1,1,1]);|
  <OrientedMatroid: 4 elements, rank 3>
  !gapprompt@gap>| !gapinput@L:=OMGeomLattice(O); LGroundSet(L);|
  <Geometric lattice: 4 atoms, rank 3>
  [ [ [ 1 ], [ 2 ], [ 3 ], [ 4 ] ], 
    [ [ 1, 2 ], [ 1, 3 ], [ 2, 3 ], 
        [ 1, 4 ], [ 2, 4 ], [ 3, 4 ] ], 
    [ [ 1, 2, 3, 4 ] ] ]
\end{Verbatim}
 }

 

\subsection{\textcolor{Chapter }{OMCocircuits (for IsOrientedMatroid)}}
\logpage{[ 6, 2, 5 ]}\nobreak
\hyperdef{L}{X83F110A97A80B140}{}
{\noindent\textcolor{FuncColor}{$\triangleright$\enspace\texttt{OMCocircuits({\mdseries\slshape OM})\index{OMCocircuits@\texttt{OMCocircuits}!for IsOrientedMatroid}
\label{OMCocircuits:for IsOrientedMatroid}
}\hfill{\scriptsize (attribute)}}\\
\textbf{\indent Returns:\ }
A list 



 Computes the (signed) cocircuits of the oriented matroid OM. 

 The result is a list of sign vectors. 

 
\begin{Verbatim}[commandchars=!@|,fontsize=\small,frame=single,label=Example]
  !gapprompt@gap>| !gapinput@O:=OrientedMatroid([[1,0],[0,1],[1,1]]); OMCocircuits(O);|
  <OrientedMatroid: 3 elements, rank 2>
  [ [ 0, 1, 1 ], [ 0, -1, -1 ], [ 1, 0, 1 ], 
    [ -1, 0, -1 ], [ -1, 1, 0 ], [ 1, -1, 0 ] ]
\end{Verbatim}
 }

 

\subsection{\textcolor{Chapter }{TopeGraph (for IsOrientedMatroid)}}
\logpage{[ 6, 2, 6 ]}\nobreak
\hyperdef{L}{X7B35E88A79675D1E}{}
{\noindent\textcolor{FuncColor}{$\triangleright$\enspace\texttt{TopeGraph({\mdseries\slshape OM})\index{TopeGraph@\texttt{TopeGraph}!for IsOrientedMatroid}
\label{TopeGraph:for IsOrientedMatroid}
}\hfill{\scriptsize (attribute)}}\\
\textbf{\indent Returns:\ }
A graph 



 Constructs the tope graph of the oriented matroid \mbox{\texttt{\mdseries\slshape OM}}. The topes are obtained as the maximal covectors from \texttt{OMCovectors(OM)}. Two topes are connected by an edge if their separating set consists of one
element. 
\begin{Verbatim}[commandchars=!@|,fontsize=\small,frame=single,label=Example]
  !gapprompt@gap>| !gapinput@O := OrientedMatroid([[1,0],[0,1],[1,1]]);|
  <OrientedMatroid: 3 elements, rank 2>
  !gapprompt@gap>| !gapinput@G := ShallowCopy(TopeGraph(O));;|
  !gapprompt@gap>| !gapinput@Vertices(G);|
  [ 1 .. 6 ]
  !gapprompt@gap>| !gapinput@IsSimpleGraph(G);|
  true
  !gapprompt@gap>| !gapinput@UndirectedEdges(G);|
  [ [ 1, 3 ], [ 1, 5 ], [ 2, 4 ], [ 2, 6 ], [ 3, 6 ], [ 4, 5 ] ]
\end{Verbatim}
 }

 

\subsection{\textcolor{Chapter }{OMCovectors (for IsOrientedMatroid)}}
\logpage{[ 6, 2, 7 ]}\nobreak
\hyperdef{L}{X7B2D2370811476C9}{}
{\noindent\textcolor{FuncColor}{$\triangleright$\enspace\texttt{OMCovectors({\mdseries\slshape OM})\index{OMCovectors@\texttt{OMCovectors}!for IsOrientedMatroid}
\label{OMCovectors:for IsOrientedMatroid}
}\hfill{\scriptsize (attribute)}}\\
\textbf{\indent Returns:\ }
A list of lists 



 Computes all covectors of the oriented matroid OM. The result is returned as a
list of lists, where the $k$\texttt{\symbol{45}}th inner list is a list of covectors corresponding to
faces of dimension $k-1$. 
\begin{Verbatim}[commandchars=!@|,fontsize=\small,frame=single,label=Example]
  !gapprompt@gap>| !gapinput@O := OrientedMatroid([[1,0],[0,1],[1,1]]);|
  <OrientedMatroid: 3 elements, rank 2>
  !gapprompt@gap>| !gapinput@OMCovectors(O);|
  [ [ [ -1, 1, -1 ], [ 1, -1, 1 ], [ -1, -1, -1 ], 
      [ 1, 1, 1 ], [ -1, 1, 1 ], [ 1, -1, -1 ] ],
    [ [ 0, 1, 1 ], [ 0, -1, -1 ], [ 1, 0, 1 ], 
      [ -1, 0, -1 ], [ -1, 1, 0 ], [ 1, -1, 0 ] ],
    [ [ 0, 0, 0 ] ] ]
\end{Verbatim}
 }

 

\subsection{\textcolor{Chapter }{SalvettiComplex (for IsOrientedMatroid)}}
\logpage{[ 6, 2, 8 ]}\nobreak
\hyperdef{L}{X7ED4D5138012EB6C}{}
{\noindent\textcolor{FuncColor}{$\triangleright$\enspace\texttt{SalvettiComplex({\mdseries\slshape OM, or, A})\index{SalvettiComplex@\texttt{SalvettiComplex}!for IsOrientedMatroid}
\label{SalvettiComplex:for IsOrientedMatroid}
}\hfill{\scriptsize (attribute)}}\\
\textbf{\indent Returns:\ }
A \texttt{FacePoset} 



 Constructs the Salvetti complex \cite{Salvetti1987_SalCpx} associated with the oriented matroid OM or the real hyperplane arrangement \mbox{\texttt{\mdseries\slshape A}}. 
\begin{Verbatim}[commandchars=!@|,fontsize=\small,frame=single,label=Example]
  !gapprompt@gap>| !gapinput@O := OrientedMatroid([[1,0],[0,1],[1,1]]);|
  <OrientedMatroid: 3 elements, rank 2>
  !gapprompt@gap>| !gapinput@SalvettiComplex(O);|
  <FacePoset of dimension 2 with f-vector [ 6, 12, 6 ]>
  !gapprompt@gap>| !gapinput@A:=HyperplaneArrangement([[1,0],[0,1],[1,1]]);|
  <HyperplaneArrangement: 3 hyperplanes in 2-space>
  !gapprompt@gap>| !gapinput@SalvettiComplex(A);|
  <FacePoset of dimension 2 with f-vector [ 6, 12, 6 ]>
\end{Verbatim}
 }

 

\subsection{\textcolor{Chapter }{SalvettiComplex (for IsHyperplaneArrangement)}}
\logpage{[ 6, 2, 9 ]}\nobreak
\hyperdef{L}{X8333113D7A7A3A91}{}
{\noindent\textcolor{FuncColor}{$\triangleright$\enspace\texttt{SalvettiComplex({\mdseries\slshape A})\index{SalvettiComplex@\texttt{SalvettiComplex}!for IsHyperplaneArrangement}
\label{SalvettiComplex:for IsHyperplaneArrangement}
}\hfill{\scriptsize (attribute)}}\\


 See \texttt{SalvettiComplex} (\ref{SalvettiComplex:for IsOrientedMatroid}). }

 

\subsection{\textcolor{Chapter }{FPGroundSet (for IsFacePoset)}}
\logpage{[ 6, 2, 10 ]}\nobreak
\hyperdef{L}{X85AB2A797ABE8967}{}
{\noindent\textcolor{FuncColor}{$\triangleright$\enspace\texttt{FPGroundSet({\mdseries\slshape FP})\index{FPGroundSet@\texttt{FPGroundSet}!for IsFacePoset}
\label{FPGroundSet:for IsFacePoset}
}\hfill{\scriptsize (attribute)}}\\
\textbf{\indent Returns:\ }
A list 



 Returns the ground set of the face poset FP. 
\begin{Verbatim}[commandchars=!@|,fontsize=\small,frame=single,label=Example]
  !gapprompt@gap>| !gapinput@S := SalvettiComplex(O);|
  <FacePoset of dimension 2 with f-vector [ 6, 12, 6 ]>
  !gapprompt@gap>| !gapinput@FPGroundSet(S);|
  [
    [ 
      [ [ -1, 1, -1 ], [ -1, 1, -1 ] ], 
      [ [ 1, -1, 1 ], [ 1, -1, 1 ] ], 
      [ [ -1, -1, -1 ], [ -1, -1, -1 ] ], 
      [ [ 1, 1, 1 ], [ 1, 1, 1 ] ], 
      [ [ -1, 1, 1 ], [ -1, 1, 1 ] ], 
      [ [ 1, -1, -1 ], [ 1, -1, -1 ] ]
    ],
    [
      [ [ -1, 0, -1 ], [ -1, 1, -1 ] ], 
      [ [ -1, 1, 0 ], [ -1, 1, -1 ] ], 
      [ [ 1, 0, 1 ], [ 1, -1, 1 ] ], 
      [ [ 1, -1, 0 ], [ 1, -1, 1 ] ], 
      [ [ 0, -1, -1 ], [ -1, -1, -1 ] ], 
      [ [ -1, 0, -1 ], [ -1, -1, -1 ] ], 
      [ [ 0, 1, 1 ], [ 1, 1, 1 ] ], 
      [ [ 1, 0, 1 ], [ 1, 1, 1 ] ], 
      [ [ 0, 1, 1 ], [ -1, 1, 1 ] ], 
      [ [ -1, 1, 0 ], [ -1, 1, 1 ] ], 
      [ [ 0, -1, -1 ], [ 1, -1, -1 ] ], 
      [ [ 1, -1, 0 ], [ 1, -1, -1 ] ]
    ],
    [
      [ [ 0, 0, 0 ], [ -1, 1, -1 ] ], 
      [ [ 0, 0, 0 ], [ 1, -1, 1 ] ], 
      [ [ 0, 0, 0 ], [ -1, -1, -1 ] ], 
      [ [ 0, 0, 0 ], [ 1, 1, 1 ] ], 
      [ [ 0, 0, 0 ], [ -1, 1, 1 ] ], 
      [ [ 0, 0, 0 ], [ 1, -1, -1 ] ]
    ]
  ]
\end{Verbatim}
 }

 

\subsection{\textcolor{Chapter }{FPOrder (for IsFacePoset)}}
\logpage{[ 6, 2, 11 ]}\nobreak
\hyperdef{L}{X7E28E057841D6ECD}{}
{\noindent\textcolor{FuncColor}{$\triangleright$\enspace\texttt{FPOrder({\mdseries\slshape FP})\index{FPOrder@\texttt{FPOrder}!for IsFacePoset}
\label{FPOrder:for IsFacePoset}
}\hfill{\scriptsize (attribute)}}\\
\textbf{\indent Returns:\ }
A function 



 Returns the order function of the face poset FP. }

 

\subsection{\textcolor{Chapter }{OMChirotope (for IsOrientedMatroid)}}
\logpage{[ 6, 2, 12 ]}\nobreak
\hyperdef{L}{X854947857F5E5D29}{}
{\noindent\textcolor{FuncColor}{$\triangleright$\enspace\texttt{OMChirotope({\mdseries\slshape OM})\index{OMChirotope@\texttt{OMChirotope}!for IsOrientedMatroid}
\label{OMChirotope:for IsOrientedMatroid}
}\hfill{\scriptsize (attribute)}}\\
\textbf{\indent Returns:\ }
A function 



 Returns the chirotope given as a function from $\binom{E}{r} \to \{0,\pm 1\}$, where $E$ is the groundset of $OM$ and $r$ is the rank of \mbox{\texttt{\mdseries\slshape OM}}. 
\begin{Verbatim}[commandchars=!@|,fontsize=\small,frame=single,label=Example]
  !gapprompt@gap>| !gapinput@O:=OrientedMatroid(2,3,[1,1,1]); c:=OMChirotope(O);|
  <OrientedMatroid: 3 elements, rank 2>
  function( S ) ... end
  !gapprompt@gap>| !gapinput@List(Combinations([1..3],2),x->c(x));|
  [ 1, 1, 1 ]
\end{Verbatim}
 }

 

\subsection{\textcolor{Chapter }{OMRankFunction (for IsOrientedMatroid)}}
\logpage{[ 6, 2, 13 ]}\nobreak
\hyperdef{L}{X7EE8B82D792A4A64}{}
{\noindent\textcolor{FuncColor}{$\triangleright$\enspace\texttt{OMRankFunction({\mdseries\slshape OM})\index{OMRankFunction@\texttt{OMRankFunction}!for IsOrientedMatroid}
\label{OMRankFunction:for IsOrientedMatroid}
}\hfill{\scriptsize (attribute)}}\\
\textbf{\indent Returns:\ }
A function 



 Returns the rank function of the oriented matroid \mbox{\texttt{\mdseries\slshape OM}} }

 }

 
\section{\textcolor{Chapter }{Properties}}\label{Chapter_Oriented_Matroids_Section_Properties}
\logpage{[ 6, 3, 0 ]}
\hyperdef{L}{X871597447BB998A1}{}
{
  

\subsection{\textcolor{Chapter }{OMIsLinear (for IsOrientedMatroid)}}
\logpage{[ 6, 3, 1 ]}\nobreak
\hyperdef{L}{X86B70496781D3884}{}
{\noindent\textcolor{FuncColor}{$\triangleright$\enspace\texttt{OMIsLinear({\mdseries\slshape OM})\index{OMIsLinear@\texttt{OMIsLinear}!for IsOrientedMatroid}
\label{OMIsLinear:for IsOrientedMatroid}
}\hfill{\scriptsize (property)}}\\
\textbf{\indent Returns:\ }
true or false 



 Determines whether the oriented matroid \mbox{\texttt{\mdseries\slshape OM}} is linear. This is the case if the internal component \texttt{OM!.lforms} is bound. 
\begin{Verbatim}[commandchars=!@|,fontsize=\small,frame=single,label=Example]
  !gapprompt@gap>| !gapinput@A := AGpql(2,2,3); OM := OrientedMatroid(A);|
  <OrientedMatroid: 6 elements, rank 3>
  !gapprompt@gap>| !gapinput@OMIsLinear(OM);|
  true
  !gapprompt@gap>| !gapinput@O := OrientedMatroid(3,4,[1,1,1,1]);|
  <OrientedMatroid: 4 elements, rank 3>
  !gapprompt@gap>| !gapinput@OMIsLinear(O);|
  false
\end{Verbatim}
 }

 

\subsection{\textcolor{Chapter }{HasSimpleSimplex (for IsOrientedMatroid)}}
\logpage{[ 6, 3, 2 ]}\nobreak
\hyperdef{L}{X7D5893807BB9785E}{}
{\noindent\textcolor{FuncColor}{$\triangleright$\enspace\texttt{HasSimpleSimplex({\mdseries\slshape OM})\index{HasSimpleSimplex@\texttt{HasSimpleSimplex}!for IsOrientedMatroid}
\label{HasSimpleSimplex:for IsOrientedMatroid}
}\hfill{\scriptsize (property)}}\\
\textbf{\indent Returns:\ }
true or false 



 Returns whether the oriented matroid \mbox{\texttt{\mdseries\slshape OM}} has a simple simplex, i.e. a simplicial Tope $T$ such that the localizations at its vertices are boolean and there is another
rank 1 covector not adjacent to $T$. If \texttt{Rank(OM)} $\geq 3$ and this is true then its Salvetti complex can not be $K(\pi,1)$. 
\begin{Verbatim}[commandchars=!@|,fontsize=\small,frame=single,label=Example]
  !gapprompt@gap>| !gapinput@O:=OrientedMatroid([[1,0,0,0],[0,1,0,0],[0,0,1,0],[0,0,0,1],[1,1,1,1]]);|
  <OrientedMatroid: 5 elements, rank 4>
  !gapprompt@gap>| !gapinput@HasSimpleSimplex(O);|
    - tope: [ -1, -1, -1, -1, -1 ]
    - vertices: [ [ 0, 0, 0, -1, -1 ], [ 0, 0, -1, 0, -1 ], 
    [ 0, -1, 0, 0, -1 ], [ -1, 0, 0, 0, -1 ] ]
  true
  
\end{Verbatim}
 }

 

\subsection{\textcolor{Chapter }{HasSimpleTriangle (for IsOrientedMatroid)}}
\logpage{[ 6, 3, 3 ]}\nobreak
\hyperdef{L}{X80F069657B2045C0}{}
{\noindent\textcolor{FuncColor}{$\triangleright$\enspace\texttt{HasSimpleTriangle({\mdseries\slshape OM})\index{HasSimpleTriangle@\texttt{HasSimpleTriangle}!for IsOrientedMatroid}
\label{HasSimpleTriangle:for IsOrientedMatroid}
}\hfill{\scriptsize (property)}}\\
\textbf{\indent Returns:\ }
true or false. 



 Checks all rank $3$ localizations for a simple simplex, i.e. a simple triangle. 
\begin{Verbatim}[commandchars=!@|,fontsize=\small,frame=single,label=Example]
  
\end{Verbatim}
 }

 }

 
\section{\textcolor{Chapter }{Operations}}\label{Chapter_Oriented_Matroids_Section_Operations}
\logpage{[ 6, 4, 0 ]}
\hyperdef{L}{X7DE8E16C7C2D387B}{}
{
  

\subsection{\textcolor{Chapter }{IsLEquiv (for IsOrientedMatroid,IsOrientedMatroid)}}
\logpage{[ 6, 4, 1 ]}\nobreak
\hyperdef{L}{X7EE74EA57A779612}{}
{\noindent\textcolor{FuncColor}{$\triangleright$\enspace\texttt{IsLEquiv({\mdseries\slshape OM1, OM2})\index{IsLEquiv@\texttt{IsLEquiv}!for IsOrientedMatroid,IsOrientedMatroid}
\label{IsLEquiv:for IsOrientedMatroid,IsOrientedMatroid}
}\hfill{\scriptsize (operation)}}\\
\textbf{\indent Returns:\ }
true or false 



 Determines if the oriented matroids \mbox{\texttt{\mdseries\slshape OM1}} and \mbox{\texttt{\mdseries\slshape OM2}} have isomorphic geometric lattices. 
\begin{Verbatim}[commandchars=!@|,fontsize=\small,frame=single,label=Example]
  !gapprompt@gap>| !gapinput@O1:=OrientedMatroid(3,4,[1,1,1,1]);|
  <OrientedMatroid: 4 elements, rank 3>
  !gapprompt@gap>| !gapinput@O2:=OrientedMatroid([[1,0,0],[0,1,0],[0,0,1],[1,1,1]]);|
  <OrientedMatroid: 4 elements, rank 3>
  !gapprompt@gap>| !gapinput@IsLEquiv(O1,O2);|
  true
  !gapprompt@gap>| !gapinput@O2:=OrientedMatroid([[1,0,0],[0,1,0],[0,0,1],[1,1,0]]);|
  <OrientedMatroid: 4 elements, rank 3>
  !gapprompt@gap>| !gapinput@IsLEquiv(O1,O2);|
  false
\end{Verbatim}
 }

 

\subsection{\textcolor{Chapter }{HasSimpleSimplexRk (for IsOrientedMatroid, IsInt)}}
\logpage{[ 6, 4, 2 ]}\nobreak
\hyperdef{L}{X81DBFBF67AC497C4}{}
{\noindent\textcolor{FuncColor}{$\triangleright$\enspace\texttt{HasSimpleSimplexRk({\mdseries\slshape OM, k})\index{HasSimpleSimplexRk@\texttt{HasSimpleSimplexRk}!for IsOrientedMatroid, IsInt}
\label{HasSimpleSimplexRk:for IsOrientedMatroid, IsInt}
}\hfill{\scriptsize (operation)}}\\
\textbf{\indent Returns:\ }
true or false 



 Returns whether the oriented matroid \mbox{\texttt{\mdseries\slshape OM}} has a localiazion of rank \mbox{\texttt{\mdseries\slshape k}} with a simple Simplex. If this is true for $k \geq 3$ then its Salvetti complex can not be $K(\pi,1)$. 

 
\begin{Verbatim}[commandchars=!@|,fontsize=\small,frame=single,label=Example]
  !gapprompt@gap>| !gapinput@O:=OrientedMatroid([[1,0,0,0],[0,1,0,0],[0,0,1,0],[0,0,0,1],[1,1,1,0]]);|
  <OrientedMatroid: 5 elements, rank 4>
  !gapprompt@gap>| !gapinput@HasSimpleSimplex(O);|
  false
  !gapprompt@gap>| !gapinput@HasSimpleSimplexRk(O,3);|
  true
\end{Verbatim}
 }

 

\subsection{\textcolor{Chapter }{OMLocalizationRk (for IsOrientedMatroid, IsList, IsInt)}}
\logpage{[ 6, 4, 3 ]}\nobreak
\hyperdef{L}{X7C0D5D2A7C058B20}{}
{\noindent\textcolor{FuncColor}{$\triangleright$\enspace\texttt{OMLocalizationRk({\mdseries\slshape OM, F, k})\index{OMLocalizationRk@\texttt{OMLocalizationRk}!for IsOrientedMatroid, IsList, IsInt}
\label{OMLocalizationRk:for IsOrientedMatroid, IsList, IsInt}
}\hfill{\scriptsize (operation)}}\\
\textbf{\indent Returns:\ }
An oriented matroid 



 Returns the rank \mbox{\texttt{\mdseries\slshape k}} localization of \mbox{\texttt{\mdseries\slshape OM}} at the flat \mbox{\texttt{\mdseries\slshape F}}. 

 
\begin{Verbatim}[commandchars=!@|,fontsize=\small,frame=single,label=Example]
  !gapprompt@gap>| !gapinput@O:=OrientedMatroid([[1,0,0,0],[0,1,0,0],[0,0,1,0],[0,0,0,1],[1,1,1,0]]);|
  <OrientedMatroid: 5 elements, rank 4>
  !gapprompt@gap>| !gapinput@L:=OMGeomLattice(O);|
  <Geometric lattice: 5 atoms, rank 4>
  !gapprompt@gap>| !gapinput@LkFlats(L)(3);|
  [ [ 1, 2, 3, 5 ], [ 1, 2, 4 ], [ 1, 3, 4 ], [ 2, 3, 4 ], [ 1, 4, 5 ], [ 2, 4, 5 ], [ 3, 4, 5 ] ]
  !gapprompt@gap>| !gapinput@OMLocalizationRk(O,[1,2,3,5],3);|
  <OrientedMatroid: 4 elements, rank 3>
  
\end{Verbatim}
 }

 }

 
\section{\textcolor{Chapter }{Global methods}}\label{Chapter_Oriented_Matroids_Section_Global_methods}
\logpage{[ 6, 5, 0 ]}
\hyperdef{L}{X80841F8D80722081}{}
{
  

\subsection{\textcolor{Chapter }{OMOperation}}
\logpage{[ 6, 5, 1 ]}\nobreak
\hyperdef{L}{X7A515BDC8582FFE6}{}
{\noindent\textcolor{FuncColor}{$\triangleright$\enspace\texttt{OMOperation({\mdseries\slshape sv1, sv2})\index{OMOperation@\texttt{OMOperation}}
\label{OMOperation}
}\hfill{\scriptsize (function)}}\\
\textbf{\indent Returns:\ }
A list 



 Performs the oriented matroid operation on two signed vectors \texttt{sv1} and \texttt{sv2}. For each position \texttt{i}, if \texttt{sv1[i]} is nonzero, it is used; otherwise \texttt{sv2[i]} is used. Returns a new signed vector representing the result of the operation. 
\begin{Verbatim}[commandchars=!@|,fontsize=\small,frame=single,label=Example]
  !gapprompt@gap>| !gapinput@CVs:=OMCovectors(O)[2];|
  [ [ 0, 1, 1 ], [ 0, -1, -1 ], [ 1, 0, 1 ], 
    [ -1, 0, -1 ], [ -1, 1, 0 ], [ 1, -1, 0 ] ]
  !gapprompt@gap>| !gapinput@sv1 := CVs[3]; sv2 := CVs[5];|
  [ 1, 0, 1 ]
  [ -1, 1, 0 ]
  !gapprompt@gap>| !gapinput@OMOperation(sv1, sv2);|
  [ 1, 1, 1 ]
  !gapprompt@gap>| !gapinput@OMOperation(sv2, sv1);|
  [ -1, 1, 1 ]
\end{Verbatim}
 }

 

\subsection{\textcolor{Chapter }{OrderCovec}}
\logpage{[ 6, 5, 2 ]}\nobreak
\hyperdef{L}{X801F6A9B7A0A53F5}{}
{\noindent\textcolor{FuncColor}{$\triangleright$\enspace\texttt{OrderCovec({\mdseries\slshape arg})\index{OrderCovec@\texttt{OrderCovec}}
\label{OrderCovec}
}\hfill{\scriptsize (function)}}\\


 Order function for sign vectors. }

 

\subsection{\textcolor{Chapter }{IsOMEquiv}}
\logpage{[ 6, 5, 3 ]}\nobreak
\hyperdef{L}{X80CB9A9C8135A314}{}
{\noindent\textcolor{FuncColor}{$\triangleright$\enspace\texttt{IsOMEquiv({\mdseries\slshape OM1, OM2})\index{IsOMEquiv@\texttt{IsOMEquiv}}
\label{IsOMEquiv}
}\hfill{\scriptsize (function)}}\\
\textbf{\indent Returns:\ }
true or false 



 Determines, if the oriented matroids \mbox{\texttt{\mdseries\slshape OM1}} and \mbox{\texttt{\mdseries\slshape OM2}} are isomorphic. }

 

\subsection{\textcolor{Chapter }{SeparatingSet}}
\logpage{[ 6, 5, 4 ]}\nobreak
\hyperdef{L}{X793FF1428136B222}{}
{\noindent\textcolor{FuncColor}{$\triangleright$\enspace\texttt{SeparatingSet({\mdseries\slshape arg})\index{SeparatingSet@\texttt{SeparatingSet}}
\label{SeparatingSet}
}\hfill{\scriptsize (function)}}\\


 The separating set of two sign vectors. }

 }

 }

   
\chapter{\textcolor{Chapter }{Special arrangements}}\label{Chapter_Special_arrangements}
\logpage{[ 7, 0, 0 ]}
\hyperdef{L}{X7C9F5AAF83133F9A}{}
{
  

 The following describes functions to construct some special arrangements. 

 
\section{\textcolor{Chapter }{Global methods}}\label{Chapter_Special_arrangements_Section_Global_methods}
\logpage{[ 7, 1, 0 ]}
\hyperdef{L}{X80841F8D80722081}{}
{
  

\subsection{\textcolor{Chapter }{AGpql}}
\logpage{[ 7, 1, 1 ]}\nobreak
\hyperdef{L}{X858205A37B7EB5FA}{}
{\noindent\textcolor{FuncColor}{$\triangleright$\enspace\texttt{AGpql({\mdseries\slshape p, q, l})\index{AGpql@\texttt{AGpql}}
\label{AGpql}
}\hfill{\scriptsize (function)}}\\
\textbf{\indent Returns:\ }
A hyperplane arrangement. 



 Constructs the reflection arrangement associated to the monomial complex
reflection group $G(p,q,l)$. The hyperplanes are defined by equations of the form 
\[ x_i = \zeta^k x_j \]
 where $\zeta$ is a primitive $p$\texttt{\symbol{45}}th root of unity. 
\begin{Verbatim}[commandchars=!@|,fontsize=\small,frame=single,label=Example]
  !gapprompt@gap>| !gapinput@A := AGpql(2,1,3);|
   <HyperplaneArrangement: 9 hyperplanes in 3-space>
  !gapprompt@gap>| !gapinput@Roots(A);|
  [ [ 1, 0, 0 ], [ 0, 1, 0 ], [ 0, 0, 1 ],
    [ 1, 1, 0 ], [ 1, -1, 0 ], [ 1, 0, 1 ],
    [ 1, 0, -1 ], [ 0, 1, 1 ], [ 0, 1, -1 ] ]
  
\end{Verbatim}
 }

 

\subsection{\textcolor{Chapter }{SsS3}}
\logpage{[ 7, 1, 2 ]}\nobreak
\hyperdef{L}{X81AD6CA97F71ED33}{}
{\noindent\textcolor{FuncColor}{$\triangleright$\enspace\texttt{SsS3({\mdseries\slshape n})\index{SsS3@\texttt{SsS3}}
\label{SsS3}
}\hfill{\scriptsize (function)}}\\
\textbf{\indent Returns:\ }
A hyperplane arrangement. 



 Generates the irreducible supersolvable simplicial arrangement of rank 3 with
n hyperplanes. n must either be even oder congruent 1 mod 4. In any case, $n\geq 6$. 
\begin{Verbatim}[commandchars=!@|,fontsize=\small,frame=single,label=Example]
  !gapprompt@gap>| !gapinput@A:=SsS3(9); Roots(A);|
  <HyperplaneArrangement: 9 hyperplanes in 3-space>
  [ [ 0, 0, 1 ], [ 0, 1, 0 ], 
    [ -1/2*E(8)+1/2*E(8)^3, 1/2*E(8)-1/2*E(8)^3, 0 ], 
    [ -1, 0, 0 ], 
    [ -1/2*E(8)+1/2*E(8)^3, -1/2*E(8)+1/2*E(8)^3, 0 ], 
    [ 1/2*E(8)-1/2*E(8)^3, 1/2*E(8)-1/2*E(8)^3, 1 ], 
    [ -1/2*E(8)+1/2*E(8)^3, 1/2*E(8)-1/2*E(8)^3, 1 ], 
    [ -1/2*E(8)+1/2*E(8)^3, -1/2*E(8)+1/2*E(8)^3, 1 ], 
    [ 1/2*E(8)-1/2*E(8)^3, -1/2*E(8)+1/2*E(8)^3, 1 ] ]
  
\end{Verbatim}
 }

 

\subsection{\textcolor{Chapter }{ConnectedSubgraphArr (for IsList)}}
\logpage{[ 7, 1, 3 ]}\nobreak
\hyperdef{L}{X82DB5C6F83FDC443}{}
{\noindent\textcolor{FuncColor}{$\triangleright$\enspace\texttt{ConnectedSubgraphArr({\mdseries\slshape Es})\index{ConnectedSubgraphArr@\texttt{ConnectedSubgraphArr}!for IsList}
\label{ConnectedSubgraphArr:for IsList}
}\hfill{\scriptsize (operation)}}\\
\textbf{\indent Returns:\ }
A hyperplane arrangement. 



 Generates the connected subgraph arrangement (see \cite{CK2022_ConnSubgraph_arXiv}) of the graph with edges \mbox{\texttt{\mdseries\slshape Es}}. 

 Let $G = (N,E)$ be an undirected graph with vertex set $N = [n]$ and edge set $E$. For $I \subseteq N$, let $G[I]$ be the induced subgraph of $G$ on the set of vertices $I$. For $I \subseteq N$, define the hyperplane 
\[H_I := \ker \sum_{i \in I}x_i.\]
 The connected subgraph arrangement consists of hyperplanes $\{H_I \mid \varnothing \ne I \subseteq N \text{ if $G[I]$ is connected} \}.$ 
\begin{Verbatim}[commandchars=!@|,fontsize=\small,frame=single,label=Example]
  !gapprompt@gap>| !gapinput@A:=ConnectedSubgraphArr([[1,2],[1,3],[2,3]]); Roots(A);|
  <HyperplaneArrangement: 7 hyperplanes in 3-space>
  [ [ 1, 1, 1 ], [ 0, 1, 1 ], [ 0, 0, 1 ], [ 0, 1, 0 ], 
    [ 1, 0, 1 ], [ 1, 0, 0 ], [ 1, 1, 0 ] ]
\end{Verbatim}
 }

 }

 }

   
\chapter{\textcolor{Chapter }{Chamber and tope posets}}\label{Chapter_Chamber_and_tope_posets}
\logpage{[ 8, 0, 0 ]}
\hyperdef{L}{X7A142C9D82C31E24}{}
{
  
\section{\textcolor{Chapter }{Operation}}\label{Chapter_Chamber_and_tope_posets_Section_Operation}
\logpage{[ 8, 1, 0 ]}
\hyperdef{L}{X835880B9800BD2D8}{}
{
  

\subsection{\textcolor{Chapter }{TopePoset (for IsOrientedMatroid, IsList)}}
\logpage{[ 8, 1, 1 ]}\nobreak
\hyperdef{L}{X804189147ADFB6A2}{}
{\noindent\textcolor{FuncColor}{$\triangleright$\enspace\texttt{TopePoset({\mdseries\slshape OM, BTope})\index{TopePoset@\texttt{TopePoset}!for IsOrientedMatroid, IsList}
\label{TopePoset:for IsOrientedMatroid, IsList}
}\hfill{\scriptsize (operation)}}\\
\textbf{\indent Returns:\ }
A TopePoset object 



 Constructs the tope poset of the oriented matroid \mbox{\texttt{\mdseries\slshape OM}} with respect to the base tope \mbox{\texttt{\mdseries\slshape BTope}} (given as sign vector). 

 
\begin{Verbatim}[commandchars=!@|,fontsize=\small,frame=single,label=Example]
  !gapprompt@gap>| !gapinput@O:=OrientedMatroid([[1,0],[0,1],[1,1]]); |
\end{Verbatim}
 }

 }

 
\section{\textcolor{Chapter }{Attributes}}\label{Chapter_Chamber_and_tope_posets_Section_Attributes}
\logpage{[ 8, 2, 0 ]}
\hyperdef{L}{X7C701DBF7BAE649A}{}
{
  

\subsection{\textcolor{Chapter }{TPGroundSet (for IsTopePoset)}}
\logpage{[ 8, 2, 1 ]}\nobreak
\hyperdef{L}{X81854C4D816A501A}{}
{\noindent\textcolor{FuncColor}{$\triangleright$\enspace\texttt{TPGroundSet({\mdseries\slshape TP})\index{TPGroundSet@\texttt{TPGroundSet}!for IsTopePoset}
\label{TPGroundSet:for IsTopePoset}
}\hfill{\scriptsize (attribute)}}\\
\textbf{\indent Returns:\ }
A list of lists. 



 The ground set of the tope poset \mbox{\texttt{\mdseries\slshape TP}}. It is a list of list, recording the separating elements of the given tope
from the base tope. }

 

\subsection{\textcolor{Chapter }{TPBaseTope (for IsTopePoset)}}
\logpage{[ 8, 2, 2 ]}\nobreak
\hyperdef{L}{X82CEA0EB7B0AF294}{}
{\noindent\textcolor{FuncColor}{$\triangleright$\enspace\texttt{TPBaseTope({\mdseries\slshape TP})\index{TPBaseTope@\texttt{TPBaseTope}!for IsTopePoset}
\label{TPBaseTope:for IsTopePoset}
}\hfill{\scriptsize (attribute)}}\\
\textbf{\indent Returns:\ }
A sign vector. 



 Returns the bast tope of the tope poset \mbox{\texttt{\mdseries\slshape TP}}. 

 }

 }

 }

 \def\bibname{References\logpage{[ "Bib", 0, 0 ]}
\hyperdef{L}{X7A6F98FD85F02BFE}{}
}

\bibliographystyle{alpha}
\bibliography{hyparr}

\addcontentsline{toc}{chapter}{References}

\def\indexname{Index\logpage{[ "Ind", 0, 0 ]}
\hyperdef{L}{X83A0356F839C696F}{}
}

\cleardoublepage
\phantomsection
\addcontentsline{toc}{chapter}{Index}


\printindex

\immediate\write\pagenrlog{["Ind", 0, 0], \arabic{page},}
\newpage
\immediate\write\pagenrlog{["End"], \arabic{page}];}
\immediate\closeout\pagenrlog
\end{document}
